% \iffalse meta-comment
%
% Parts of this file, Copyright (C) 1994 by Jerry Leichter
% Copyright (C) 2016-2024 by Pieter van Oostrum <pieter@vanoostrum.org>
% -------------------------------------------------------
%
% This file may be distributed and/or modified under the
% conditions of the LaTeX Project Public License, either version 1.3
% of this license or (at your option) any later version.
% The latest version of this license is in:
%
%    http://www.latex-project.org/lppl.txt
%
% and version 1.3 or later is part of all distributions of LaTeX
% version 2005/12/01 or later.
%
% \fi
%
% \iffalse
%<*driver>
\ProvidesFile{\jobname.dtx}
%</driver>
% This is the variable part of the preamble
%
%<multirow|bigstrut>%% Copyright (C) 1994 by Jerry Leichter
%<bigdelim>%% Copyright (C) 1994 by \O ystein Bache
%% Copyright (C) 2016-2024 by Pieter van Oostrum <pieter@vanoostrum.org>
%%
%
%    \begin{macrocode}
%<multirow|bigstrut|bigdelim>\NeedsTeXFormat{LaTeX2e}
%<multirow>\ProvidesPackage{multirow}%
%<bigstrut>\ProvidesPackage{bigstrut}
%<bigdelim>\ProvidesPackage{bigdelim}
%<multirow|bigstrut|bigdelim>           [2024/11/12 v2.9
%<multirow>                  Span multiple rows of a table]%
%<bigstrut>                  Provide larger struts in tabulars]
%<bigdelim>                  Create big delimiters in tabular or array]
%    \end{macrocode}
%
%<*driver>
\documentclass[a4paper]{ltxdoc}
\usepackage[fontset=fandol]{ctex}
\usepackage{multirow,bigstrut,bigdelim,colortbl}
\usepackage[table]{xcolor}
\usepackage{tabulary}
\usepackage{longtable}
\usepackage{booktabs}
\usepackage{calc}
\usepackage{url}
\usepackage[colorlinks, urlcolor=blue, linktocpage]{hyperref}
\usepackage{tikz}
\usepackage{hhline}
\usepackage{etoolbox}
\AtBeginEnvironment{macrocode}{\color{red}}
\newcommand{\minitab}[2][l]{\begin{tabular}{#1}#2\end{tabular}}
\DisableCrossrefs
\CodelineIndex
\RecordChanges
\IndexPrologue{\section*{索引}%
  \markboth{索引}{索引}%
  斜体数字给出条目所在说明的页码；带下划线的数字给出宏定义所在的代码行号；罗马体数字给出条目被使用的代码行号。}
\GlossaryPrologue{\section*{变更历史}%
  \markboth{变更历史}{变更历史}}
\def\generalname{一般}
\newcommand{\env}[1]{\texttt{#1}}
\newcommand{\Package}[1]{\textsf{#1}}
\def\bsbs{\cs{\char`\\}}
\begin{document}
  \DeleteShortVerb{\|}
  \DocInput{\jobname.dtx}
\end{document}
%</driver>
% \fi
%
% \CheckSum{0}
%
% \CharacterTable
%  {Upper-case    \A\B\C\D\E\F\G\H\I\J\K\L\M\N\O\P\Q\R\S\T\U\V\W\X\Y\Z
%   Lower-case    \a\b\c\d\e\f\g\h\i\j\k\l\m\n\o\p\q\r\s\t\u\v\w\x\y\z
%   Digits        \0\1\2\3\4\5\6\7\8\9
%   Exclamation   \!     Double quote  \"     Hash (number) \#
%   Dollar        \$     Percent       \%     Ampersand     \&
%   Acute accent  \'     Left paren    \(     Right paren   \)
%   Asterisk      \*     Plus          \+     Comma         \,
%   Minus         \-     Point         \.     Solidus       \/
%   Colon         \:     Semicolon     \;     Less than     \<
%   Equals        \=     Greater than  \>     Question mark \?
%   Commercial at \@     Left bracket  \[     Backslash     \\
%   Right bracket \]     Circumflex    \^     Underscore    \_
%   Grave accent  \`     Left brace    \{     Vertical bar  \|
%   Right brace   \}     Tilde         \~}
%
% \changes{v1.7}{2016/09/13}{统一所有文件的版本号}
% \changes{v2.0}{2016/09/27}{发布 v2.0}
% \changes{v2.3}{2018/08/03}{小缺陷修复}
% \changes{v2.4}{2018/12/30}{在 \texttt{bigstrut.sty} 中添加 \leavevmode}
% \changes{v2.4}{2019/01/01}{使 \cs{multirow} 中的 \meta{width} 与 \meta{vmove}
% 与 \Package{calc} 兼容}
% \changes{v2.5}{2019/05/31}{解决与 \Package{colortbl} 宏包的冲突}
%
% \GetFileInfo{multirow.sty}
%
% \DoNotIndex{\#,\$,\%,\&,\@,\\,\{,\},\^,\_,\~,\ }
% \DoNotIndex{\@ne}
% \DoNotIndex{\advance,\begingroup,\catcode,\closein}
% \DoNotIndex{\closeout,\day,\def,\edef,\else,\empty,\endgroup}
%
% \title{\Package{multirow}、
%   \Package{bigstrut} 与
%   \Package{bigdelim} 宏包}
% \author{Pieter van Oostrum\thanks{在册的“现行作者”}\\
%   \O ystein Bache\\
%   Jerry Leichter\thanks{文档最初由 Robin
%     Fairbairns 整理}}
%
% \date{发布于 \filedate，版本 \fileversion}
%
% \maketitle
%
% \tableofcontents
%
% \section{引言}
% 这些宏包为标准的 \LaTeX{}
% \env{tabular} 环境提供了一系列扩展。它们各自的功能如下：
% \begin{description}
% \item[\Package{multirow}] 提供一种构造，用于生成跨越表格多行的单元格；
%
% \item[\Package{bigstrut}] 生成支撑条，用于（略微）撑开所在的表格行；
%
% \item[\Package{bigdelim}] 生成尺寸合适的定界符（如花括号、圆括号或方括号），
%   放入单个跨行单元中，用以标示其他各行之间的关系；
%
% \end{description}
%
% \section{2.x 版本的变更}
% \label{sec:changes-2}
%
%\subsection*{版本 2.9}
%
% \begin{itemize}
% \item 文档订正：此前对 \meta{nrows} 取负值时含义的描述有误。
%
% \item 使 \texttt{multirow.sty} 与 \texttt{supertabular} 4.2a 或更高版本兼容。
%
% \end{itemize}
%
%\subsection*{版本 2.8}
%
%\begin{itemize}
%\item 为 \cs{ldelim} 与 \cs{rdelim} 增加可选参数 \meta{vmove}。
%\end{itemize}
%
%\subsection*{版本 2.7}
%
%\begin{itemize}
%\item  将 \cs{@xmultirow} 与 \cs{multirow@setcolwidth} 改为 \cs{long} 定义，以支持多段文本。
%
% 参见 \url{https://github.com/pietvo/multirow/issues/1}
%\end{itemize}
%
%\subsection*{版本 2.6}
%
% \begin{itemize}
% \item 使 \texttt{supertabular} 选项与较新版本的 \Package{supertabular} 宏包兼容。
%
% \item 当 \cs{@arstrutbox} 未定义时对其初始化，以便在 \env{array} 或 \env{tabular} 之外也能使用大型定界符。
%
% \end{itemize}
%
%\subsection*{版本 2.5a}
%
% \begin{itemize}
% \item 更新联系方式
% \end{itemize}
%
%\subsection*{版本 2.5}
%
% \begin{itemize}
% \item 解决使用 \texttt{longtable} 选项时与 \Package{colortbl} 宏包的冲突；此前两个宏包都会重定义 \cs{@cline}。
%
%
% \end{itemize}
%
%\subsection*{版本 2.4}
%
% \begin{itemize}
% \item 在 \Package{bigstrut} 中添加 \cs{leavevmode}，以强制进入水平模式。
% \item 使 \cs{multirow} 中的 \meta{width} 与 \meta{vmove} 可与 \Package{calc} 兼容。
% \end{itemize}
%
%\subsection*{版本 2.3}
%
% \begin{itemize}
% \item 在 \texttt{bigdelim.sty} 中，将大括号旁文本所用的 \cs{textrm} 替换为 \cs{textnormal}。
% \end{itemize}
%
%\subsection*{版本 2.2}
%
% \begin{itemize}
% \item 支持 \meta{nrows} 取小数值。
% \end{itemize}
%
%\subsection*{版本 2.0}
%
% \begin{itemize}
% \item \cs{multirow} 现在带有第一个可选参数 \oarg{vpos}。
% \item \meta{width} 参数可指定为 \texttt{=}，表示使用 \cs{multirow} 所在列的既定宽度。
%
% \item 为 \meta{bigstruts} 参数增加可选前缀字母（\texttt{t}、\texttt{b}）（参见第~\ref{sect-refine-bigstrut} 节）。
%
% \item 宏包选项 \texttt{debug}。
% \item 宏包选项 \texttt{longtable}，用于绕过 \Package{longtable} 的一处缺陷。参见第~\ref{sec:use-with-longtable} 节。
%
% \item 宏包选项 \texttt{supertabular}，以更好地支持 \Package{supertabular}。参见
% 第~\ref{sec:use-with-supertabular} 节。
% \item 改进了某些情形下的定位效果。
% \item 补充了大量文档。
% \item 发布版本改以 \texttt{.dtx} 文件为基础组织。
% \item 与 v1.6 向后兼容。
% \end{itemize}
%
% \section{使用 \Package{multirow}}\label{sec:multirow}
%
% \DescribeMacro{\multirow}
% \cs{multirow} 在 \env{tabular} 或
% 类似环境中放置一段文本，使其跨越多行。我们把该文本所跨越的行列区块称为跨行区块。通常它只覆盖一列，但与 \cs{multicolumn} 结合使用时可以覆盖更多列。
%
%
%
% \changes{v2.1}{2016/10/06}{依照《LaTeX 伴侣》，将文档中的 \meta{fixup} 更名为 \meta{vmove}。}
%
%
% 基本语法如下：
% \begin{quote}
%   \cmd{\multirow}\oarg{vpos}\marg{nrows}\oarg{bigstruts}\marg{width}\oarg{vmove}\marg{text}
% \end{quote}
% 其中：
% \begin{description}
% \item[\meta{vpos}] 设定文本在跨行区块中的垂直位置。
% 默认为 \texttt{[c]}，即文本垂直居中；其他选项有顶部对齐的 \texttt{[t]} 和底部对齐的 \texttt{[b]}。
%
% \item[\meta{nrows}] 为跨越的行数。本列的其余各行应留空，否则 \cs{multirow} 生成的内容会覆盖其上。
%
% 当 \meta{nrows} 取正值时，跨越的行包括当前行及其下方（\meta{nrows}-1）行；取负值时，包括当前行及其上方（-\meta{nrows}-1）行。\meta{nrows} 允许取小数，以便进行微调。
%
%
%
%
%
% \item[\meta{bigstruts}] 主要在使用了 \Package{bigstrut} 宏包时使用。此时它是所跨越各行中 \cs{bigstrut} 的使用总次数：每个 \cs{bigstrut} 计 2 次，每个 \cs{bigstrut}\oarg{x}（\meta{x} 为 \texttt{t} 或 \texttt{b}）计 1 次。默认为 0。
%
%
%
%
%
% \meta{bigstruts} 参数前还可加前缀字母 \texttt{t}、\texttt{b} 或 \texttt{tb}，以便更精细地控制。详见第~\ref{sect-refine-bigstrut} 节。前缀字母与数字之间可以有一个空格。
%
%
%
%
% \item[\meta{width}] 为文本的排版宽度。
% 特殊取值：\texttt{*} 表示采用文本参数的自然宽度；\texttt{=} 表示采用放置该 \cs{multirow} 条目的列的既定宽度。
%
%
%
% \item[\meta{vmove}] 是用于微调的长度：文本将在其原本位置之上（之下）再上移（下移，当 \meta{vmove} 为负值时）该长度。
%
%
%
% \item[\meta{text}] 为该构造的实际文本。若显式指定了宽度，文本将排在该宽度的 \cs{parbox} 中；可使用 \bsbs{} 在任意位置强制换行。
%
%
%
% 若宽度取 \texttt{*}，文本将在 LR（水平）模式下排版。此时若需要多行条目，可在文本参数中使用 \env{tabular} 或 \env{array} 环境，例如下面的 \env{minitab}。
%
%
%
%
% 当 \cs{multirow} 条目位于具有既定宽度的列中时，宽度也可取 \texttt{=}，例如 \texttt{p\{\}} 列、\env{tabularx} 中的 \texttt{X} 列，或 \env{tabulary} 中的 \texttt{L}、\texttt{C}、\texttt{R}、\texttt{J} 列。文本将排在该宽度的 \cs{parbox} 中。
%
%
%
%
% 在其他情形下使用“\texttt{=}”会得到奇怪的结果（通常是列宽过大）。
%
%
% \end{description}
%
% 从 \Package{multirow} 2.4 版起，若已加载 \Package{calc} 宏包，\meta{width} 与 \meta{vmove} 参数可写成 \Package{calc} 表达式。是否载入 \Package{calc} 由文档作者负责，\Package{multirow} 不会自动载入。
%
%
%
%
%
% 注意：\cs{multirow} 可用于 \env{tabular} 及其大多数派生环境，例如 \env{tabularx}、\env{tabulary}、\env{supertabular}、\env{ltablex}、\env{xtab}、\env{longtable}、\env{tabu}、\env{booktabs} 和 \env{ctable}。其中某些环境在特定情形下需要特别留意，详见下文。
%
%
%
%
%
%
% \DescribeMacro{\multirowsetup}
% 在展开 \meta{text} 之前，会先展开 \cs{multirowsetup} 宏，用于设置所需的特殊环境。初始时 \cs{multirowsetup} 仅包含 \cs{raggedright}，可用 \cs{renewcommand} 重新定义。
%
%
%
%
% 若要在同一单元格中同时使用 \cs{multirow} 与 \cs{multicolumn}，必须把 \cs{multirow} 放在 \cs{multicolumn} 内部；反过来则行不通。例如：
% \index{multicolumn=\verb!*+\multicolumn+|usage}
%
%
%
% \begin{verbatim}
%    \multicolumn{2}{c}{\multirow{3}{*}{多行多列}}
% \end{verbatim}
%
% \subsection{宏包选项}
% \label{sec:package-options}
%
% \DescribeMacro{multirowdebugtrue}
% \DescribeMacro{multirowdebugfalse}
% 定义了以下选项：
% \begin{description}
% \item[\texttt{debug}]
% 该选项会把跨行盒子的相关信息写入日志文件，由 \TeX{} 的 \cs{showbox} 命令完成。注意：即使没有真正的错误，这也会使 \LaTeX{} 认为编译失败。
%
%
%
% 该选项也可在文档任意位置用 \cs{multirowdebugtrue} 命令开启、用 \cs{multirowdebugfalse} 命令关闭。若此类命令紧跟在某个 \cs{multirow} 之前，则仅对该 \cs{multirow} 条目生效。
%
%
%
% \item[longtable]
% \texttt{longtable} 选项重定义 \cs{cline} 宏，以绕过 \Package{longtable} 宏包中的一处缺陷。参见第~\ref{sec:use-with-longtable} 节。
%
% \end{description}
%
%\subsection{示例}
%\label{sec:examples}
%
% 同时使用 \Package{multirow} 与 \Package{bigstrut} 的例子：
%
% \begin{quote}
% \begin{verbatim}
% \newcommand{\minitab}[2][l]{\begin{tabular}{#1}#2\end{tabular}}
% \begin{tabular}{|c|c|}
% \hline
% \multirow{4}{1in}{普通文本 g} & 列 g2a\\
%       & 列 g2b \\
%       & 列 g2c \\
%       & 列 g2d \\
% \hline
% \multirow{3}[6]*{普通文本 g} & 列 g2a\bigstrut\\\cline{2-2}
%       & 列 g2b \bigstrut\\\cline{2-2}
%       & 列 g2c \bigstrut\\
% \hline
% \multirow{4}[8]{1in}{普通文本 g，但稍长一些。} &
%                                    列 g2a\bigstrut\\\cline{2-2}
%       & 列 g2b \bigstrut\\\cline{2-2}
%       & 列 g2c \bigstrut\\\cline{2-2}
%       & 列 g2d \bigstrut\\
% \hline
% \multirow{4}*{\minitab[c]{普通\\ g 文本}} & 列 g2a\\
%       & 列 g2b \\
%       & 列 g2c \\
%       & 列 g2d \\
% \hline
% \end{tabular}
% \end{verbatim}
% \end{quote}
% 排版结果如下：
% \begin{quote}
%   \begin{tabular}{|c|c|}
%     \hline
%     \multirow{4}{1in}{普通文本 g} & 列 g2a\\
%       & 列 g2b \\
%       & 列 g2c \\
%       & 列 g2d \\
%     \hline
%     \multirow{3}[6]*{普通文本 g} & 列 g2a\bigstrut\\\cline{2-2}
%       & 列 g2b \bigstrut\\\cline{2-2}
%       & 列 g2c \bigstrut\\
%     \hline
%     \multirow{4}[8]{1in}{普通文本 g，但稍长一些。} &
%                                        列 g2a\bigstrut\\\cline{2-2}
%       & 列 g2b \bigstrut\\\cline{2-2}
%       & 列 g2c \bigstrut\\\cline{2-2}
%       & 列 g2d \bigstrut\\
%     \hline
%     \multirow{4}*{\minitab[c]{普通\\ g 文本}} & 列 g2a\\
%       & 列 g2b \\
%       & 列 g2c \\
%       & 列 g2d \\
%     \hline
%   \end{tabular}
%   \quad
%   \begin{tabular}{ll}
%     \multirow{4}{*}{普通情形} & \\
%       & \\
%       & \\
%       & \\
%     \multirow{3}[6]*{使用 \cs{bigstrut}，\meta{width} 取 \texttt{*}} & \bigstrut\\
%       & \bigstrut\\
%       & \bigstrut\\
%     \multirow{4}[8]{*}{使用 \cs{bigstrut}，\meta{width} 取常规值} & \bigstrut\\
%       & \bigstrut\\
%       & \bigstrut\\
%       & \bigstrut\\
%     \multirow{4}*{\cs{multirow} 中的多行文本} & \\
%       & \\
%       & \\
%       & \\
%   \end{tabular}
%
% \end{quote}
% \vspace{1ex}
% 下面是在 \env{tabulary} 中使用 \texttt{=} 宽度说明符的例子（注意：\texttt{=} 两侧的花括号可以省略）：
%
%
% \begin{quote}
% \index{extrarowheight=\verb!*+\extrarowheight+|usage}
% \begin{verbatim}
% \setlength{\extrarowheight}{2pt}
% \begin{tabulary}{11cm}{|L|L|L|}
%   \hline
%   人人生而自由，在尊严和权利上一律平等。 &
%      \multirow{2}{=}{人人有资格享受本宣言所载的一切权利和
%      自由，不分种族、肤色、性别、语言、宗教、
%      政治或其他见解、国籍或社会出身、财产、
%      出生或其他身份等任何区别。
%      }
%   & 人人有权享有生命、自由和人身安全。 \\
% \cline{1-1}\cline{3-3}
%      任何人不得加以奴役；奴隶制和奴隶买卖
%      均应予以禁止。 &
%   & 任何人不得加以酷刑，或施以残忍、不人道
%   或侮辱性的待遇或刑罚。 \\
% \hline
% \end{tabulary}
% \end{verbatim}
% \end{quote}
%
% \begin{quote}
% \index{extrarowheight=\verb!*+\extrarowheight+|usage}
% \setlength{\extrarowheight}{2pt}
% \begin{tabulary}{11cm}{|L|L|L|}
%   \hline
%   人人生而自由，在尊严和权利上一律平等。 &
%      \multirow{2}{=}{人人有资格享受本宣言所载的一切权利和
%      自由，不分种族、肤色、性别、语言、宗教、
%      政治或其他见解、国籍或社会出身、财产、
%      出生或其他身份等任何区别。}
%
%   & 人人有权享有生命、自由和人身安全。 \\
% \cline{1-1}\cline{3-3}
%      任何人不得加以奴役；奴隶制和奴隶买卖
%      均应予以禁止。 &
%   & 任何人不得加以酷刑，或施以残忍、不人道
%   或侮辱性的待遇或刑罚。 \\
% \hline
% \end{tabulary}
% \end{quote}
%
% 关于本例，有几点需要说明：
% \begin{enumerate}
% \item  中间一列是 \cs{multirow} 条目。你可能以为它会垂直居中，但实际上并非如此。这是因为 \cs{multirow} 并不知道盒子的高度；它对高度所能做的唯一估计是行数${}\times{}$单行常规高度。它试图在该空间内居中，但本例中该空间偏低，因此文本位于盒子顶部。若要居中，需要给出 \meta{vmove} 参数把它向下平移。
%
%
%
%
%
%
%
% \item 我们使用了 2pt 的 \cs{extrarowheight}\footnote{这只有在载入 \Package{array} 宏包时才可用，而 \Package{tabulary} 会自动载入它。}，以在 \cs{hline} 与后续文本之间留出一些间距。但该间距并不作用于 \cs{multirow}，因为它被假定为居中。此时可把 \meta{vpos} 参数取为 \texttt{[t]}，\cs{multirow} 便会作出恰当的定位。
%
%
%
%
%
%
% \end{enumerate}
%
% 下面改用取负值的 \meta{nrows}。
%
% \begin{quote}
% \index{extrarowheight=\verb!*+\extrarowheight+|usage}
% \begin{verbatim}
% \setlength{\extrarowheight}{2pt}
% \begin{tabulary}{11cm}{|L|L|L|}
%   \hline
%   人人生而自由，在尊严和权利上一律平等。 &
%   & 人人有权享有生命、自由和人身安全。 \\
% \cline{1-1}\cline{3-3}
%      任何人不得加以奴役；奴隶制和奴隶买卖
%      均应予以禁止。 &
%      \multirow{-2}{=}[12mm]{人人有资格享受本宣言所载的一切权利和
%      自由，不分种族、肤色、性别、语言、宗教、
%      政治或其他见解、国籍或社会出身、财产、
%      出生或其他身份等任何区别。
%      }
%   & 任何人不得加以酷刑，或施以残忍、不人道
%   或侮辱性的待遇或刑罚。 \\
% \hline
% \end{tabulary}
% \end{verbatim}
% \end{quote}
%
% \begin{quote}
% \index{extrarowheight=\verb!*+\extrarowheight+|usage}
% \setlength{\extrarowheight}{2pt}
% \begin{tabulary}{11cm}{|L|L|L|}
%   \hline
%   人人生而自由，在尊严和权利上一律平等。 &
%   & 人人有权享有生命、自由和人身安全。 \\
% \cline{1-1}\cline{3-3}
%      任何人不得加以奴役；奴隶制和奴隶买卖
%      均应予以禁止。 &
%      \multirow{-2}{=}[12mm]{人人有资格享受本宣言所载的一切权利和
%      自由，不分种族、肤色、性别、语言、宗教、
%      政治或其他见解、国籍或社会出身、财产、
%      出生或其他身份等任何区别。}
%
%   & 任何人不得加以酷刑，或施以残忍、不人道
%   或侮辱性的待遇或刑罚。 \\
% \hline
% \end{tabulary}
% \end{quote}
%
% 本例中文本会在底部一行附近居中，从而向下超出盒子。因此我们施加了 12mm 的 \meta{vmove}。\meta{vmove} 通常需要反复试验才能确定合适值。
%
%
%
% \subsection{微调}
% \label{sec:fine-tuning}
%
% 若所跨越的某一行异常高大，或者使用了 \Package{bigstrut} 宏包且 \cs{bigstrut} 在跨越各行的中心线两侧分布不对称，垂直居中可能不准确。此时应使用 \meta{vmove} 参数。有时直接给 \meta{nrows} 取更大的值（包括小数）会更有效。示例参见第~\ref{sec:dealing-with-tall} 节。
%
%
%
%
%
%
% 要完全处理好下伸字母几乎不可能。文本会先按居中排版，但其基线可能与旁边内容的基线不一致，尤其是旁边内容有下伸字母而 \meta{text} 没有时。这会造成轻微的错位。所能做的，至多是利用 \meta{vmove} 可选参数对 \meta{text} 作最后的微调。
%
%
%
%
%
% （提示：若采用 \texttt{.1ex} 这样的度量单位，那么在改变字号后，\meta{vmove} 很可能仍然适用。）
%
%
%
% \cs{multirow} 主要是为 \env{tabular}（而非 \env{array}）环境设计的。若某些行中有大型公式，它在 \env{array} 环境中可能效果不佳；此时可用 \meta{vmove} 参数细化结果。
%
%
%
%
% 某些情形下，你可能希望跨行条目与其他行单元格的顶部对齐，例如其中有一个大号大写字母时。当使用 \meta{vpos} = \texttt{[t]} 时，对齐的是基线，而这在本例中并不合适。此时可用 \meta{vmove} 参数完成定位，让 \LaTeX{} 计算平移量。例如：
%
%
%
%
%
%
% \begin{quote}
% \begin{verbatim}
% \usepackage{calc}
% \newlength{\shiftdown}
% \setlength{\shiftdown}{\heightof{\Huge\bfseries B}-\heightof{f}}
% . . .
% \begin{tabular}{cll}
%   \toprule
%   \multirow[t]{5}{*}[-\shiftdown]{\Huge\bfseries B}
%                     & 甲  & 示例文本一    \\
%                     & 乙  & 示例文本二    \\
%                     & 丙  & 示例文本三  \\
%                     & 丁  & 示例文本四 \\
%   \bottomrule
% \end{tabular}
% \end{verbatim}
% \end{quote}
%
% \begin{quote}
% \newlength{\shiftdown}
% \setlength{\shiftdown}{\heightof{\Huge\bfseries B}-\heightof{f}}
% \begin{tabular}{cll}
%   \toprule
%   \multirow[t]{5}{*}[-\shiftdown]{\Huge\bfseries B}
%                     & 甲  & 示例文本一    \\
%                     & 乙  & 示例文本二       \\
%                     & 丙  & 示例文本三  \\
%                     & 丁  & 示例文本四 \\
%   \bottomrule
% \end{tabular}
% \end{quote}
% 从 \Package{multirow} 2.4 版起，\meta{vmove} 参数也可直接使用表达式 \verb|-\heightof{\Huge|-\verb|\bfseries B}|-\verb|\heightof{f}|，以代替 \verb|-\shiftdown|。
%
%
% \subsection{跨行条目与着色单元格}\label{sec:mult-color-cells}
%
% 若将 \cs{multirow} 与 \Package{colortbl} 宏包（或带 \texttt{[table]} 选项的 \Package{xcolor} 宏包）配合使用，而又想给含 \cs{multirow} 的列着色，就需要采取一些预防措施。\Package{colortbl} 的工作方式是逐个单元格着色。因此当 \meta{nrows} 取正值时，\Package{colortbl} 会先给顶部单元格着色，随后 \cs{multirow} 从该单元格开始排版 \meta{nrows} 个单元格，之后 \Package{colortbl} 再给其余单元格着色，从而把该区域的文字遮盖掉。解决办法是把 \cs{multirow} 放在最后一行，并取负的 \meta{nrows} 值。例如：
% \index{colortbl=\Package{colortbl}|usage}
%
% \index{xcolor=\Package{xcolor}|usage}
%
%
%
%
%
%
%
%
%
% \begin{quote}
% \begin{verbatim}
% \begin{tabular}{l>{\columncolor{yellow}}l}
%   甲甲 & \\
%   乙乙 & \\
%   丙丙 & \multirow{-3}*{丁丁}\\
% \end{tabular}
% \end{verbatim}
% \end{quote}
% 排版结果如下：
% \begin{quote}
%   \begin{tabular}{l>{\columncolor{yellow}}l}
%     甲甲 & \\
%     乙乙 & \\
%     丙丙 & \multirow{-3}*{丁丁}\\
%   \end{tabular}
% \end{quote}
%
% \index{hhline=\Package{hhline}|usage}
% 当着色的跨行单元格与 \Package{hhline} 宏包配合使用时，着色的跨行单元格中可能出现一些白色条纹。例如：
%
%
% \begin{quote}
% \begin{verbatim}
%   \begin{tabular}{|>{\columncolor{red}}c|c|}
%   \hline
%   \bfseries 第一列 & \bfseries 第二列\\ \hline
%   数据一 & 932\\ \hline
%    & 239\\ \hhline{|~|-|}
%    & 137\\ \hhline{|~|-|}
%   \multirow{-3}{*}{更多数据} & 319\\ \hline
%   末行数据 & 132\\ \hline
% \end{tabular}
% \end{verbatim}
% \end{quote}
%
% \begin{quote}
%   \begin{tabular}{|>{\columncolor{red}}c|c|}
%   \hline
%   \bfseries 第一列 & \bfseries 第二列\\ \hline
%   数据一 & 932\\ \hline
%    & 239\\ \hhline{|~|-|}
%    & 137\\ \hhline{|~|-|}
%   \multirow{-3}{*}{更多数据} & 319\\ \hline
%   末行数据 & 132\\ \hline
% \end{tabular}
% \end{quote}
%
% 解决办法是在着色的跨行单元格中，放置颜色相同的水平横线。
%
% \begin{quote}
% \begin{verbatim}
%   \begin{tabular}{|>{\columncolor{red}}c|c|}
%   \hline
%   \bfseries 第一列 & \bfseries 第二列\\ \hline
%   数据一 & 932\\ \hline
%    & 239\\ \hhline{|>{\arrayrulecolor{red}}->{\arrayrulecolor{black}}|-|}
%    & 137\\ \hhline{|>{\arrayrulecolor{red}}->{\arrayrulecolor{black}}|-|}
%   \multirow{-3}{*}{更多数据} & 319\\ \hline
%   末行数据 & 132\\ \hline
% \end{tabular}
% \end{verbatim}
% \end{quote}
%
% \begin{quote}
%   \begin{tabular}{|>{\columncolor{red}}c|c|}
%   \hline
%   \bfseries 第一列 & \bfseries 第二列\\ \hline
%   数据一 & 932\\ \hline
%    & 239\\ \hhline{|>{\arrayrulecolor{red}}->{\arrayrulecolor{black}}|-|}
%    & 137\\ \hhline{|>{\arrayrulecolor{red}}->{\arrayrulecolor{black}}|-|}
%   \multirow{-3}{*}{更多数据} & 319\\ \hline
%   末行数据 & 132\\ \hline
% \end{tabular}
% \end{quote}
%
% \subsection{细化 \meta{bigstruts} 参数}
% \label{sect-refine-bigstrut}
%
% 只要各行没有“异常”高度，\cs{multirow} 就能根据 \meta{nrows} 与 \meta{bigstruts} 算出所需跨行盒子的高度。但有些情况下这不足以正确定位盒子，尤其是首行顶部和（或）末行底部含有 \cs{bigstrut} 时。此时应向 \cs{multirow} 提供额外信息：在 \meta{bigstruts} 参数前加一个（或两个）字母，指明这两种情况中存在哪一种。
%
%
%
%
%
%
%
%
% 请看下面的例子：\\（这些例子中我们设置了 \verb+\setlength{\bigstrutjot}{10pt}+，以使效果清晰可见）\\[5mm]
%
% \begin{minipage}[c]{0.23\linewidth}
%   \setlength{\bigstrutjot}{10pt}
%   \begin{tabular}{|c|c|}
%     \hline
%     \multirow{3}[1]{*}{跨行} & T \bigstrut[t] \\
%     \cline{2-2}
%     &  X \\
%     \cline{2-2}
%     &  B \\
%     \hline
%   \end{tabular} \\[2mm]
%   \begin{tabular}{|c|c|}
%     \hline
%     \multirow{3}[t 1]{*}{跨行} & T \bigstrut[t] \\
%     \cline{2-2}
%     &  X \\
%     \cline{2-2}
%     &  B \\
%     \hline
%   \end{tabular}
% \end{minipage}
% \begin{minipage}[c]{0.75\linewidth}
% \begin{verbatim}
%   \begin{tabular}{|c|c|}
%     \hline
%     \multirow{3}[1]{*}{Multirow} & T \bigstrut[t] \\
%     \cline{2-2}
%     &  X \\
%     \cline{2-2}
%     &  B \\
%     \hline
%   \end{tabular}
% \end{verbatim}
% \end{minipage}
% \\[5mm]
% 上例上方盒子中，文本“跨行”本应居中，但由于首行有 \cs{bigstrut[t]}，它略偏下。可以把 \meta{bigstruts} 参数取为“t 1”，表示顶部有一个支撑条来纠正。下方盒子即采用此法，使用了 \verb+\multirow{3}[t 1]{*}{Multirow}+。
%
%
%
%
%
% 第二个例子：\\[5mm]
% \begin{minipage}{0.23\linewidth}
%   \setlength{\bigstrutjot}{10pt}
%   \begin{tabular}{|c|c|}
%     \hline
%     & T \\
%     \cline{2-2}
%     &  X \\
%     \cline{2-2}
%     \multirow[t]{-3}[1]{*}{跨行} &  B \bigstrut[b] \\
%     \hline
%   \end{tabular} \\[2mm]
%   \begin{tabular}{|c|c|}
%     \hline
%     & T \\
%     \cline{2-2}
%     &  X \\
%     \cline{2-2}
%     \multirow[t]{-3}[b 1]{*}{跨行} &  B \bigstrut[b] \\
%     \hline
%   \end{tabular}
% \end{minipage}
% \begin{minipage}{0.75\linewidth}
% \begin{verbatim}
%   \begin{tabular}{|c|c|}
%     \hline
%     & T \\
%     \cline{2-2}
%     &  X \\
%     \cline{2-2}
%     \multirow[t]{-3}[1]{*}{Multirow} &  B \bigstrut[b] \\
%     \hline
%   \end{tabular}
% \end{verbatim}
% \end{minipage}
% \\[5mm]
% 上方盒子中，\verb+\multirow[t]+ 本应与 \texttt{T} 处于同一高度，但它偏高，因为底部有一个 \cs{bigstrut}。可把 \meta{bigstruts} 参数指定为“b 1”来纠正，即使用 \verb+\multirow[t]{-3}[b 1]{*}{Multirow}+。
%
%
%
%
% \vspace{5mm}
% 前缀的可选取值如下：
% \begin{description}
% \item[\texttt{t}] 顶部有支撑条，即首行中有 \cs{bigstrut} 或 \cs{bigstrut[t]}。
%
% \item[\texttt{b}] 底部有支撑条，即末行中有 \cs{bigstrut} 或 \cs{bigstrut[b]}。
%
% \item[\texttt{tb}] 两者都有。注意：不能写成 \texttt{bt}。
%
% \end{description}
% 字母与数字之间的空格可有可无。
% 请注意，前缀与 \cs{multirow} 位于首行还是末行无关。
%
%
% \subsection{与 \Package{longtable} 配合使用}
% \label{sec:use-with-longtable}
%
% \index{longtable=\env{longtable}|usage}
% \index{longtabu=\env{longtabu}|usage}
% 可以在 \env{longtable}（及其派生环境 \env{longtabu}）中使用 \cs{multirow}。但要注意：当跨行条目靠近页面底部时，应避免 longtable 在其中途分页。例如：
%
%
%
%
% \noindent
% \begin{minipage}[t]{0.4\linewidth}
% \begin{longtable}{|l|l|l|}
% \ldots & \ldots & \ldots \\
% \hline
% 9 月 21 日 & 09:00 & 事件 1 \\
% \hline
% 9 月 22 日 & 10:00 & 事件 2 \\
% \hline
% \multirow{3}*{9 月 23 日} & 10:00 & 事件 3 \\
%                         & 12:00 & 事件 4 \\
%                         & 15:00 & 事件 5 \\
% \hline
% 9 月 24 日 & 09:00 & 事件 6 \\
% \hline
% \ldots & \ldots & \ldots \\
% \end{longtable}
% \end{minipage}
% \begin{minipage}[t]{0.58\linewidth}
% \begin{verbatim}
% \begin{longtable}{|l|l|l|}
% \ldots & \ldots & \ldots \\
% \hline
% 9 月 21 日 & 09:00 & 事件 1 \\
% \hline
% 9 月 22 日 & 10:00 & 事件 2 \\
% \hline
% \multirow{3}*{9 月 23 日} & 10:00 & 事件 3 \\
%                         & 12:00 & 事件 4 \\
%                         & 15:00 & 事件 5 \\
% \hline
% 9 月 24 日 & 09:00 & 事件 6 \\
% \hline
% \ldots & \ldots & \ldots \\
% \end{longtable}
% \end{verbatim}
% \end{minipage}
%
% 本例中，若“9 月 23 日”条目靠近页面底部，应阻止分页发生在该条目中部。做法是把中间各行以 \verb+\\*+ 结尾，而不是 \verb+\\+。
%
%
%
% \begin{verbatim}
% \multirow{3}*{9 月 23 日} & 10:00 & 事件 3 \\*
%                         & 12:00 & 事件 4 \\*
%                         & 15:00 & 事件 5 \\
% \end{verbatim}
%
% 但 \Package{longtable} 存在一处缺陷：当 \verb+\\*+ 后面紧跟 \cs{cline} 时，它会失效，如下例所示。
%
%
%
% \begin{verbatim}
% \multirow{3}*{9 月 23 日} & 10:00 & 事件 3 \\*
% \cline{2-3}
%                         & 12:00 & 事件 4 \\*
% \cline{2-3}
%                         & 15:00 & 事件 5 \\
% \end{verbatim}
%
% \Package{multirow} 提供了 \texttt{longtable} 宏包选项，它重定义 \cs{cline}，使 \verb+\\*+ 在其后紧跟 \cs{cline} 时也能生效。这段代码来自 David Carlisle。
%
%
%
% \subsection{与 \Package{supertabular} 配合使用}
% \label{sec:use-with-supertabular}
%
% \index{supertabular=\env{supertabular}|usage}
% \index{xtab=\env{xtab}|usage}
% 使用 \Package{supertabular} 宏包（或其增强版 \Package{xtab}）时，同样需要在分页时把跨行的各行保持在一起。遗憾的是，\Package{supertabular} 没有指定“禁止分页”的方式，即 \verb+\\*+ 并不能阻止分页。
%
%
%
%
%
% 因此 \Package{multirow} 提供了 \texttt{supertabular} 宏包选项，它重定义 \env{supertabular} 内部的 \verb+\\*+ 以阻止分页。应当用它来结束跨行区块中的中间各行。但与 \env{longtable} 不同，这并不会使 \env{supertabular} 考虑在 \cs{multirow} 之前分页，因此表格可能变得过长。
%
%
%
%
%
%
%
% \DescribeMacro{\STneed}
% 因此，在指定 \texttt{supertabular} 选项时，\Package{multirow} 还提供了 \cs{STneed} 命令，用于在 \env{supertabular} 中说明本页还需要多少空间；若空间不足，便会在该处分页。例如：
%
%
%
%
% \begin{verbatim}
% \tabletail{\hline}
% \begin{supertabular}{|l|l|l|}
% \ldots & \ldots & \ldots \\
% \hline
% 9 月 20 日 & 10:00 & 事件 1\\
% \hline
% 9 月 21 日 & 10:00 & 事件 2\\
% \hline
% 9 月 22 日 & 10:00 & 事件 3\\
% \hline
% \STneed {2cm}
% \multirow{3}*{9 月 23 日} & 10:00 & 事件 4\\*
% \cline{2-3}
%                         & 12:00 & 事件 5 \\*
% \cline{2-3}
%                         & 15:00 & 事件 6 \\
% \hline
%  9 月 24 日 & 09:00 & 事件 7 \\
% \hline
% \ldots & \ldots & \ldots \\
% \end{verbatim}
%
% 该取值需要反复试验才能确定合适值；有时把所需空间适当夸大一些反而更好。
%
%
% \subsection{处理过高的条目}
% \label{sec:dealing-with-tall}
%
% 有时某些行会比预期更高。本节给出处理这类情形的一些思路，分两种情况：
%
%
% \begin{enumerate}
% \item 第一种情况：跨行区块之外存在异常高大的行，此时 \cs{multirow} 的定位可能出错。这是因为 \cs{multirow} 没有各行高度的信息。例如某单元格中输入了大型公式，或在 \texttt{[\{\}]} 列中输入了多行段落。示例：
%
%
%
%
%
%
% \begin{quote}
% \begin{verbatim}
% \begin{tabular}{| l | l | p{4cm} |}
%   \hline
%   \multirow{3}*{第 38 周} & 周一 & 全天多雨\\
%   \cline{2-3}
%                     & 周二 & 晴间多云\\
%   \cline{2-3}
%                     & 周三 & 晴朗多日照；
%                     但大风会使
%                     气温下降。 \\
%   \hline
% \end{tabular}
% \end{verbatim}
%  \end{quote}
% \begin{quote}
% \begin{tabular}{| l | l | p{4cm} |}
%   \hline
%   \multirow{3}*{第 38 周} & 周一 & 全天多雨\\
%   \cline{2-3}
%                     & 周二 & 晴间多云\\
%   \cline{2-3}
%                     & 周三 & 晴朗多日照；
%                     但大风会使
%                     气温下降。 \\
%   \hline
% \end{tabular}
% \end{quote}
% 该 \cs{multirow} 位于第二行，因为它指定跨越 3 行。但本例中第二行并非垂直居中位置，因为第三行要高得多。
%
%
%
% 要纠正这一点，可以使用 \meta{vmove} 参数。但本例中更简单的做法是假定 \cs{multirow} 跨越 6 行（最后一列的总行数），即用 \verb+\multirow{6}+\ldots{}，得到：
%
%
%
% \begin{quote}
% \begin{tabular}{| l | l | p{4cm} |}
%   \hline
%   \multirow{6}*{第 38 周} & 周一 & 全天多雨\\
%   \cline{2-3}
%                     & 周二 & 晴间多云\\
%   \cline{2-3}
%                     & 周三 & 晴朗多日照；
%                     但大风会使
%                     气温下降。 \\
%   \hline
% \end{tabular}
% \end{quote}
% \item 第二种情况：\cs{multirow} 条目比周围常规行更高。此时跨行文本会伸出其区块之外。必须把其他行加高，而这是 \cs{multirow} 无法自动完成的。
%
%
%
%
% 示例：
% （请勿将其视为医疗建议；药名均为虚构。）
% \begin{verbatim}
% \begin{tabular}{| p{2mm} l | p{5cm} |}
%  \hline
%  \multicolumn{2}{|l|}{\textbf{药物与剂量}}
%                               & \textbf{可能的副作用} \\
%  \hline
%  \multicolumn{2}{|l|}{斯皮诺}
%                      & \multirow{3}={意识模糊、
%                        幻觉、呼吸急促、
%                        癫痫（惊厥）；
%                        胃部不适、烧心；严重恶心、
%                        呕吐、胃痛或轻度头痛。} \\
%  \cline{1-2}
%   & 起始：200 毫克/日 & \\
%  \cline{1-2}
%    & 维持：100–400 毫克/日 & \\
%  \hline
%  \multicolumn{2}{|l|}{康扎克}
%                      & \multirow{3}={焦虑；紧张；
%                       失眠；食欲减退；轻度心动过缓；
%                       窦房结传导减慢；体重下降；
%                       光敏；低钠血症；
%                       性功能障碍（男女均可发生）；可能
%                       影响糖尿病患者的血糖控制。} \\
%  \cline{1-2}
%   & 起始：10 毫克/日 & \\
%  \cline{1-2}
%    & 维持：10–40 毫克/日 & \\
%  \hline
% \end{tabular}
% \end{verbatim}
% \begin{tabular}{| p{2mm} l | p{5cm} |}
%  \hline
%  \multicolumn{2}{|l|}{\textbf{药物与剂量}}
%                               & \textbf{可能的副作用} \\
%  \hline
%  \multicolumn{2}{|l|}{斯皮诺}
%                      & \multirow{3}={意识模糊、
%                        幻觉、呼吸急促、
%                        癫痫（惊厥）；
%                        胃部不适、烧心；严重恶心、
%                        呕吐、胃痛或轻度头痛。} \\
%  \cline{1-2}
%   & 起始：200 毫克/日 & \\
%  \cline{1-2}
%    & 维持：100–400 毫克/日 & \\
%  \hline
%  \multicolumn{2}{|l|}{康扎克}
%                      & \multirow{3}={焦虑；紧张；
%                       失眠；食欲减退；轻度心动过缓；
%                       窦房结传导减慢；体重下降；
%                       光敏；低钠血症；
%                       性功能障碍（男女均可发生）；可能
%                       影响糖尿病患者的血糖控制。} \\
%  \cline{1-2}
%   & 起始：10 毫克/日 & \\
%  \cline{1-2}
%    & 维持：10–40 毫克/日 & \\
%  \hline
% \end{tabular}
% \vspace{2cm}
% 两个 \cs{multirow} 条目都过高：第一个伸入第二个条目区域，第二个则伸出表格之外。
%
%
% 有两种纠正办法：
% 最简单的办法是添加额外的空行来覆盖重叠空间。第一个条目需补 2 行，第二个需补 4 行。因此在第三条 \verb+\hline+ 之前插入两行 \verb+ & & \\+，在最后一条 \verb+\hline+ 之前插入四行。这样表格就恰好合适：
%
%
%
%
% \vspace{5mm}
% \begin{tabular}{| p{2mm} l | p{5cm} |}
%   \hline
%   \multicolumn{2}{|l|}{\textbf{药物与剂量}}
%                               & \textbf{可能的副作用} \\
%   \hline
%  \multicolumn{2}{|l|}{斯皮诺}
%                      & \multirow{3}={意识模糊、
%                        幻觉、呼吸急促、
%                        癫痫（惊厥）；
%                        胃部不适、烧心；严重恶心、
%                        呕吐、胃痛或轻度头痛。} \\
%   \cline{1-2}
%     & 起始：200 毫克/日 & \\
%   \cline{1-2}
%      & 维持：100–400 毫克/日 & \\
%     & & \\
%     & & \\
%   \hline
%  \multicolumn{2}{|l|}{康扎克}
%                      & \multirow{3}={焦虑；紧张；
%                       失眠；食欲减退；轻度心动过缓；
%                       窦房结传导减慢；体重下降；
%                       光敏；低钠血症；
%                       性功能障碍（男女均可发生）；可能
%                       影响糖尿病患者的血糖控制。} \\
%   \cline{1-2}
%     & 起始：10 毫克/日 & \\
%   \cline{1-2}
%      & 维持：10–40 毫克/日 & \\
%     & & \\
%     & & \\
%     & & \\
%     & & \\
%   \hline
% \end{tabular}
%
% 第二种办法是在垂直方向上撑高常规行，使其与跨行条目相配。本表字号为 10pt，每行总高 12pt。第一个条目需额外 24pt（2 行）；由于该空间要分摊到 3 行，每行 8pt，使行总高达到 20pt。
%
%
%
%
%
% 常规行高 8.4pt、深 3.6pt（合计 12pt）。我们可以在上、下各加 4pt，或任何合计为 8pt 的组合。此处我取高 12pt、深 8pt，用宽度为 0 的 \cs{rule} 实现。
%
%
%
% \DescribeMacro{\mystrut}
% \verb+\newcommand{\mystrut}{\rule[-8pt]{0pt}{20pt}}+
% 并把 \cs{mystrut} 放入前 3 行的每一行。自定义支撑条可以完全掌控版式：例如可以让某些行多占空间，或把全部空间都放在最后一行。
%
%
%
%
% 第二个条目需额外 48pt（4 行）。我们在每行使用 \cs{bigstrut}，即每行 16pt；由于一个 \cs{bigstrut} 等于 2\cs{bigstrutjot}，把 \cs{bigstrutjot} 设为 8pt。
%
%
% \Package{booktabs} 会在横线上下增加一些额外垂直空间，因此使用常规 \env{tabular} 环境时支撑条宜略大，而使用 \Package{booktabs} 时宜略小。经试验，\cs{bigstrutjot} 取 \texttt{7pt} 即足够。当然，第二个跨行条目还加上了 \texttt{[tb6]} 的 \meta{bigstruts} 参数。注意：用自定义支撑条无法做到这一点，除非“作弊”。
%
%
%
%
%
%
%
%
% \index{booktabs=\env{booktabs}|usage}
% 改用 \Package{booktabs} 后，代码如下：
% \begin{verbatim}
% \newcommand{\mystrut}{\rule[-8pt]{0pt}{20pt}}
% \setlength{\bigstrutjot}{7pt}
% \begin{tabular}{ p{2mm} l  p{5cm} }
%  \toprule
%  \multicolumn{2}{l}{\textbf{药物与剂量}}
%                               & \textbf{可能的副作用} \\
%  \cmidrule(r){1-2} \cmidrule(l){3-3}
%  \multicolumn{2}{l}{斯皮诺} \mystrut
%                      & \multirow{3}={意识模糊、
%                        幻觉、呼吸急促、
%                        癫痫（惊厥）；
%                        胃部不适、烧心；严重恶心、
%                        呕吐、胃痛或轻度头痛。} \\
%  \cmidrule(r){1-2}
%     & 起始：200 毫克/日 \mystrut & \\
%  \cmidrule(r){1-2}
%      & 维持：100–400 毫克/日 \mystrut & \\
%  \midrule
%  \multicolumn{2}{l}{康扎克} \bigstrut
%                      & \multirow{3}[tb6]={焦虑；紧张；
%                       失眠；食欲减退；轻度心动过缓；
%                       窦房结传导减慢；体重下降；
%                       光敏；低钠血症；
%                       性功能障碍（男女均可发生）；可能
%                       影响糖尿病患者的血糖控制。} \\
%  \cmidrule(r){1-2}
%    & 起始：10 毫克/日 \bigstrut & \\
%  \cmidrule(r){1-2}
%    & 维持：10–40 毫克/日 \bigstrut  & \\
%  \bottomrule
% \end{tabular}
% \end{verbatim}
%
% \newcommand{\mystrut}{\rule[-8pt]{0pt}{20pt}}
% \setlength{\bigstrutjot}{7pt}
% \begin{tabular}{ p{2mm} l  p{5cm} }
%  \toprule
%  \multicolumn{2}{l}{\textbf{药物与剂量}}
%                               & \textbf{可能的副作用} \\
%  \cmidrule(r){1-2} \cmidrule(l){3-3}
%  \multicolumn{2}{l}{斯皮诺} \mystrut
%                      & \multirow{3}={意识模糊、
%                        幻觉、呼吸急促、
%                        癫痫（惊厥）；
%                        胃部不适、烧心；严重恶心、
%                        呕吐、胃痛或轻度头痛。} \\
%  \cmidrule(r){1-2}
%     & 起始：200 毫克/日 \mystrut & \\
%  \cmidrule(r){1-2}
%      & 维持：100–400 毫克/日 \mystrut & \\
%  \midrule
%  \multicolumn{2}{l}{康扎克} \bigstrut
%                      & \multirow{3}[tb6]={焦虑；紧张；
%                       失眠；食欲减退；轻度心动过缓；
%                       窦房结传导减慢；体重下降；
%                       光敏；低钠血症；
%                       性功能障碍（男女均可发生）；可能
%                       影响糖尿病患者的血糖控制。} \\
%  \cmidrule(r){1-2}
%    & 起始：10 毫克/日 \bigstrut & \\
%  \cmidrule(r){1-2}
%    & 维持：10–40 毫克/日 \bigstrut  & \\
%  \bottomrule
% \end{tabular}
% \end{enumerate}
%
% \section{使用 \Package{bigstrut}}
%
% \DescribeMacro{\bigstrut}
% \DescribeMacro{\bigstrutjot}
% \cs{bigstrut}\oarg{x} 生成一个支撑条（宽度为 0 的横线），比标准 \env{array}/\env{tabular} 支撑条高出、低出或同时高出低出 \cs{bigstrutjot}（默认为 \texttt{2pt}）。把它用在紧邻 \cs{hline} 的表格单元格中，可留出额外一点间距——正如《TeXbook》（第 246 页）所说：“这是一处提升带框表格观感的小细节，可视为质量的标志。”
%
%
%
%
%
%
% 虽然也可以在 \env{array} 中使用 \cs{bigstrut}，但通常意义不大，因为 \env{array} 本来就通过 \cs{jot} 拉开了间距。
%
%
% \cs{bigstrut[t]} 增加高度；\cs{bigstrut[b]} 增加深度；单独的 \cs{bigstrut} 两者都增加。用法：在 \cs{hline} 紧接\emph{之后}的行用 \cs{bigstrut[t]}；在其紧接\emph{之前}的行用 \cs{bigstrut[b]}；前后都有 \cs{hline} 时用 \cs{bigstrut}。
%
%
%
%
%
% \cs{bigstrut} 之后的空格会被忽略，即使它带有可选参数；其之前的空格一般也会被忽略（通过一次 \cs{unskip}）。
%
%
%
% \DescribeMacro{\bigstrutjot}
% 注意：\Package{multirow} 宏包也会使用 \cs{bigstrutjot}。两个宏包同时使用时，先后顺序不限，因为它们都会检查对方是否已定义 \cs{bigstrutjot}。但二者设定的默认值不同：仅使用 \Package{multirow} 时，\cs{bigstrutjot} 设为 \texttt{3pt}；使用 \Package{bigstrut} 时（无论是否同时用 \Package{multirow}），\cs{bigstrutjot} 为 \texttt{2pt}。
%
%
%
%
%
%
%
% \section{使用 \Package{bigdelim}}
%
% 该宏包用于在同时使用 \Package{multirow} 的 \env{tabular} 或 \env{array} 环境中工作。
%
%
% \DescribeMacro{\ldelim}
% \DescribeMacro{\rdelim}
% 用法语法为：
% \begin{quote}
%   \cs{ldelim} ( \marg{n}\oarg{vmove}\marg{width}\oarg{text}\\
%   \cs{rdelim} ) \marg{n}\oarg{vmove}\marg{width}\oarg{text}
% \end{quote}
% 这两个命令用在 \env{tabular} 或 \env{array} 的某一列中；它们生成一个大型圆括号、花括号或其他定界符，从所在行起跨越 \meta{n} 行。后续各行对应的单元格必须显式留空。
%
%
%
%
%
% 第一个参数是所用的定界符，例如 \cs{\char`\{}
% \cs{\char`\}} \texttt{[} \texttt{]} \texttt{(} \texttt{)}\,---\,实际上，凡是能与 \cs{left} 或 \cs{right} 配合使用的符号均可。
%
%
% 示例如下：
% \begin{quote}
% \begin{verbatim}
% \begin{equation}
%   \begin{array}{ccccccc}
%     \ldelim({4}{4mm} & x & x & x & x & \rdelim){4}{4mm}  \\
%     & x & x & x & x &  & i \\
%     & x & x & x & x &  & j \\
%     & x & x & x & x &  \\
%     &   & u & v &   &
%   \end{array}
% \end{equation}
% \end{verbatim}
% \end{quote}
%
% \begin{quote}
% \begin{equation}
%   \begin{array}{ccccccc}
%     \ldelim({4}{4mm} & x & x & x & x & \rdelim){4}{4mm}  \\
%     & x & x & x & x &  & i \\
%     & x & x & x & x &  & j \\
%     & x & x & x & x &  \\
%     &   & u & v &   &
%   \end{array}
% \end{equation}
% \end{quote}
%
% 可选参数 \meta{vmove} 是用于微调的长度：定界符（连同可选的 \meta{text}）将在其原本位置之上（之下）再上移（下移，当 \meta{vmove} 为负值时）该长度。这与 \cs{multirow} 相同，但注意此处 \meta{vmove} 位于 \meta{width} 之前。
%
%
%
%
%
% 使用 \cs{ldelim} 时，可选 \meta{text} 排在 \cs{ldelim} 左侧并居中；使用 \cs{rdelim} 时，则排在 \cs{rdelim} 右侧。\meta{width} 参数是为定界符及其文本预留的空间；与 \Package{multirow} 一样，\meta{width} 可取 \texttt{*}。对比如下：
%
%
%
%
%
% \begin{quote}
% \begin{minipage}{0.6\linewidth}
% \begin{verbatim}
% \begin{tabular}{p{2em}l}
%   \ldelim\{{3}{*}[类型] & dvi \\
%                         & ps \\
%                         & pdf \\
% \end{tabular}
% \end{verbatim}
% \end{minipage}
% \begin{minipage}{0.35\linewidth}
% \begin{tabular}{p{2em}l}
%   \ldelim\{{3}{*}[类型] & dvi \\
%                             & ps \\
%                             & pdf \\
%    \end{tabular}
%    \\[2ex]
% \end{minipage}
%
% \begin{minipage}{0.6\linewidth}
% \begin{verbatim}
% \begin{tabular}{l@{\,}l}
%   \ldelim\{{3}{*}[类型] & dvi \\
%                         & ps \\
%                         & pdf \\
% \end{tabular}
% \end{verbatim}
% \end{minipage}
% \begin{minipage}{0.35\linewidth}
%   \begin{tabular}{l@{\,}l}
%   \ldelim\{{3}{*}[类型] & dvi \\
%                          & ps \\
%                          & pdf \\
%   \end{tabular}
%   \\[2ex]
% \end{minipage}
% \end{quote}
% 第一个例子中我们“作弊”了：通过使用过窄的列宽，吞掉了一部分列间距，代价是出现“Overfull \cs{hbox}”警告。第二个例子采用正规做法：插入 \verb+@{\,}+，用一个窄间距替换默认列间距。
%
%
%
%
%
% 这些命令也可以在扩展区块的最后一行使用，并取负的 \meta{n} 参数。这在配合 \Package{colortbl} 或 \Package{xcolor} 宏包时很有用（参见第~\ref{sec:multirow} 节中关于 \Package{multirow} 的讨论）。若存在异常高大的行，可能需要增大 \meta{n}（可取小数）。
%
% \index{colortbl=\Package{colortbl}|usage}
%
% \index{xcolor=\Package{xcolor}|usage}
%
%
%
%
% 若有与花括号相交的横线，建议使用 \Package{hhline} 宏包来绘制这些线。
% \index{hhline=\Package{hhline}|usage}
%
%
% 若用 \verb|@{}| 缩小或消除列间距，并配合 \Package{colortbl} 或 \Package{xcolor} 宏包使用着色背景（命令 \verb|\columncolor|、\verb|\rowcolor|、\verb|\rowcolors| 或 \verb|\cellcolor|），你会发现花括号的一部分被切掉。实际上它是被相邻单元格的颜色覆盖了。示例如下：
%
%
%
% \index{columncolor=\verb!*+\columncolor+|usage}
% \index{rowcolor=\verb!*+\rowcolor+|usage}
% \index{rowcolors=\verb!*+\rowcolors+|usage}
% \index{cellcolor=\verb!*+\cellcolor+|usage}
%
%
%
% \begin{quote}
% \begin{minipage}{0.6\linewidth}
% \begin{verbatim}
% \rowcolors{2}{green!25}{green!75}
% \begin{tabular}{ c @{}c c c }
%                  & 1 & 2 & 3 \\
%                  & 4 & 5 & 6 \\
% \ldelim\{{-3}{*} & 7 & 8 & 9 \\
% \end{tabular}
% \end{verbatim}
% \end{minipage}
% \begin{minipage}{0.35\linewidth}
% \rowcolors{2}{green!25}{green!75}
% \begin{tabular}{ c @{}c c c }
%                  & 1 & 2 & 3 \\
%                  & 4 & 5 & 6 \\
% \ldelim\{{-3}{*} & 7 & 8 & 9 \\
% \end{tabular}
%    \\[2ex]
% \end{minipage}
% \end{quote}
%
% 这并非 \Package{multirow} 或 \Package{bigdelim} 的问题；即使 \verb|@{}| 之前那一列是普通文本，也会出现同样情况。原因是这些着色命令会把颜色扩展到两侧的列间距上，以免出现着色间隙。这些所谓\textit{外延}的大小，每侧为 \verb|\tabcolsep|（使用 \env{array} 时为 \verb|\arraycolsep|）。但使用 \verb|@{}| 时不存在这样的列间距，扩展部分就会覆盖前一列的一部分。
%
%
%
%
%
%
%
% 解决办法是在表格列头中用 \verb|\columncolor| 命令把左侧\textit{外延}显式设为 0pt，如 \verb|>{\columncolor{white}[0pt][\tabcolsep]}|。遗憾的是，显式指定颜色 \texttt{white} 会取消该列的透明度，因此该列中未显式着色的单元格也会受影响。若 \env{tabular} 背景为白色，通常不易察觉；但如果背景颜色被改变（例如用 \verb|\pagecolor| 命令），就应使用该背景色而非 \texttt{white}。遗憾的是，目前没有只指定\textit{外延}而不同时指定颜色的命令。
%
%
%
%
%
%
%
%
%
%
%
% 下面的例子即采用了这一做法。为使表格列头简洁，我们把该定义放入 \verb|\newcolumn| 命令（使用 \Package{array} 宏包）。
%
%
%
% \begin{quote}
% \begin{minipage}{0.6\linewidth}
% \begin{verbatim}
% \newcolumntype{z}{@{}>{\columncolor{white}[0pt][\tabcolsep]}}
% \rowcolors{2}{green!25}{green!75}
% \begin{tabular}{ c zc c c }
%                  & 1 & 2 & 3 \\
%                  & 4 & 5 & 6 \\
% \ldelim\{{-3}{*} & 7 & 8 & 9 \\
% \end{tabular}
% \end{verbatim}
% \end{minipage}
% \begin{minipage}{0.35\linewidth}
% \newcolumntype{z}{@{}>{\columncolor{white}[0pt][\tabcolsep]}}
% \rowcolors{2}{green!25}{green!75}
% \begin{tabular}{ c zc c c }
%                  & 1 & 2 & 3 \\
%                  & 4 & 5 & 6 \\
% \ldelim\{{-3}{*} & 7 & 8 & 9 \\
% \end{tabular}
%    \\[2ex]
% \end{minipage}
% \end{quote}
%
% 若想把段落型文本作为可选参数，可以把它放入 \cs{parbox}。另一种办法是增加一个额外列，把文本放在该列的 \cs{multirow} 中，如下所示：
%
% \begin{quote}
% \begin{verbatim}
% \begin{tabular}{l@{}l@{}l}
%    dvi   & \rdelim\}{3}{1em} & \multirow{3}{4cm}{这些是常用的 \TeX 输出格式，
%         } \\
%    ps    & & \\
%    pdf   & & \\
% \end{tabular}
% \end{verbatim}
% \end{quote}
% \begin{quote}
% \begin{tabular}{l@{}l@{}l}
%    dvi   & \rdelim\}{3}{1em} & \multirow{3}{4cm}{这些是常用的 \TeX{} 输出格式，
%         } \\
%    ps   & & \\
%    pdf   & & \\
% \end{tabular}
% \end{quote}
% 注意，这里使用 \texttt{@\{\}} 消除列间距，使文本紧贴花括号。
%
% \section{联系方式}
%
% Pieter van Oostrum\\
% 邮箱：\texttt{pieter@vanoostrum.org}\\
% 主页：\url{http://pieter.vanoostrum.org}
% \\[1ex]
% 源代码位于 Github：\\
% \url{https://github.com/pietvo/multirow}\\
% 缺陷报告请提交至：\\
% \url{https://github.com/pietvo/multirow/issues}
%
% \StopEventually{%
% \PrintChanges
% \PrintIndex}
%
% \section{实现细节}
%
%
%\subsection{\Package{multirow} 宏包}
%\label{sec:pack-multirow}
% \iffalse
%<*multirow>
% \fi
%
% \changes{multirow v1.0}{}{匿名发布，基于一则 Usenet 帖子}
% \changes{multirow v1.1}{}{使其可在无 \texttt{bigstrut.sty} 的情况下工作（Pieter van Oostrum）}
% \changes{multirow v1.2}{}{由 Jerry Leichter 为同一目的修改，但采用
%   不同的方法，以与 \texttt{bigstrut.sty} 正确配合}
% \changes{multirow v1.2a}{}{由 Pieter van Oostrum 修改：定位时改用 \cs{vskip}
%   而非 \cs{raise}，避免在调整量较大时把行撑得过高}
%
% \changes{multirow v1.3}{}{由 Pieter van Oostrum 修改，使其在 p{} 列中正常工作
%    （添加了 \cs{leavevmode}）}
% \changes{multirow v1.4}{}{由 Pieter van Oostrum 修改：检查宽度取 * 的特殊情形。此时使用文本参数的自然宽度，该参数在 LR 模式下处理。}
%
%
%
% \changes{multirow v1.5}{}{由 Pieter van Oostrum 修改：在 \texttt{\cs{hbox}{\#5}\cs{vfill}} 之后添加了 \texttt{\%}。\\
%       在 \#5 两侧添加 \cs{strut} 以改善垂直定位。
%       增加了 \meta{nrows} 取负值时的处理代码。}
% \changes{multirow v1.6}{2004/05/05}{由 Pieter van Oostrum 修改：在 \texttt{\cs{advance}\cs{multirow@dima}\#4} 之后把一个空格替换为 \cs{relax}}
%
%
% \changes{v1.9a}{2016/09/23}{实现 \texttt{debug} 选项。}
%
% \begin{macro}{\ifmultirowdebug}
% \begin{macro}{\multirowdebugtrue}
% \begin{macro}{\multirowdebugfalse}
% 这是一个用于开启/关闭调试（显示生成盒子的内容）的布尔开关。由 \texttt{debug} 宏包选项开启，\cs{newif} 初始将其置为假（关闭）。
%
%
%    \begin{macrocode}
\newif\ifmultirowdebug
\DeclareOption{debug}{\multirowdebugtrue}
%    \end{macrocode}
% \end{macro}
% \end{macro}
% \end{macro}
%
% \changes{v1.9b}{2016/09/24}{实现 \texttt{longtable} 选项。}
% \changes{v2.5}{2019/05/31}{使 \cs{@cline} 的重定义与 \Package{colortbl} 宏包兼容。}
%
% \begin{macro}{\cline}
% \texttt{longtable} 宏包选项重定义 \cs{cline} 宏，以绕过 \Package{longtable} 中的一处缺陷。参见第~\ref{sec:use-with-longtable} 节\footnote{感谢 David Carlisle，见\href{http://tex.stackexchange.com/questions/52100/longtable-multirow-problem-with-cline-and-nopagebreak\#answer-52101}{他在 stackexchange 上的回答}。}。
%
%
%
% 首先检查宏 \cs{CT@arc} 是否已定义。若已定义，说明已载入 \Package{colortbl} 宏包。由于 \Package{colortbl} 也会重定义 \cs{@cline}，我们在自身重定义 \cs{@cline} 时必须把这一点考虑在内。
%
%
%
%    \begin{macrocode}
\DeclareOption{longtable}{%
\AtBeginDocument{%
\@ifundefined{CT@arc}
{\def\@cline#1-#2\@nil{%
  \omit
  \@multicnt#1%
  \advance\@multispan\m@ne
  \ifnum\@multicnt=\@ne\@firstofone{&\omit}\fi
  \@multicnt#2%
  \advance\@multicnt-#1%
  \advance\@multispan\@ne
  \leaders\hrule\@height\arrayrulewidth\hfill
  \cr
  \noalign{\nobreak\vskip-\arrayrulewidth}}}
{\def\@cline#1-#2\@nil{%
  \omit
  \@multicnt#1%
  \advance\@multispan\m@ne
  \ifnum\@multicnt=\@ne\@firstofone{&\omit}\fi
  \@multicnt#2%
  \advance\@multicnt-#1%
  \advance\@multispan\@ne
  {\CT@arc@\leaders\hrule\@height\arrayrulewidth\hfill}%
  \cr
  \noalign{\nobreak\vskip-\arrayrulewidth}}}
}}
%    \end{macrocode}
% \end{macro}
%
% \changes{v1.9b}{2016/09/26}{实现 \texttt{supertabular} 选项与 \cs{STneed} 命令。}
%
% \texttt{supertabular} 宏包选项重定义 \env{supertabular} 内部的 \verb+\\*+。该重定义延迟到 \verb+\begin{document}+ 时才执行。
%
%
%
% \texttt{Supertabular} 4.1f 及更高版本需要调用 \cs{ST@save@lineno} 才能正常工作，但更早版本中不存在该命令，无法使用。因此 \cs{ST@tabularcr} 与 \cs{MRST@cr} 的定义会根据 \cs{ST@save@lineno} 是否定义而有所不同。此外还有一些其他差异，某些中间版本的 \texttt{supertabular} 可能需要更细致的适配，目前暂处理到这一步。
%
%
%
%
%
%
%    \begin{macrocode}
\DeclareOption{supertabular}{%
\AtBeginDocument{%
%    \end{macrocode}
% \begin{macro}{\ST@tabularcr}
% \changes{v2.6}{2020/12/26}{调整定义，以兼容现代版本的 \texttt{supertabular}}
%
% 该宏是 \env{supertabular} 内部 \verb+\\+ 的定义。我们检查是否带有 \verb+*+：若有则调用自己的版本，否则调用 \env{supertabular} 的版本。先给出 4.1f 之前旧版 \texttt{supertabular} 所用的版本，再给出新版。
%
%
%
%
%    \begin{macrocode}
\ifx\ST@save@lineno\undefined
\def\ST@tabularcr{%
  {\ifnum0=`}\fi
  \@ifstar{\MRST@xtabularcr}{\ST@xtabularcr}}
%    \end{macrocode}
% \end{macro}
% \begin{macro}{\MRST@cr}
% \cs{MRST@cr} 是 \cs{ST@cr} 的精简副本。它完成 \env{supertabular} 所占空间的全部簿记工作，但不负责分页部分。
%
%
% \changes{v2.9}{2024/11/1}{调整定义，以兼容 \texttt{supertabular} 4.2a 或更高版本。}
%
%    \begin{macrocode}
\def\MRST@cr{%
  \noalign{%
    \ifnum\ST@pboxht<\ST@lineht
      \global\advance\ST@pageleft -\ST@lineht
      \global\ST@prevht\ST@lineht
    \else
      \global\advance\ST@pageleft -\ST@pboxht
      \global\advance\ST@pageleft -0.1\ST@pboxht
      \global\advance\ST@pageleft -\ST@stretchht
      \global\ST@prevht\ST@pboxht
      \global\ST@pboxht\z@
    \fi
    \global\advance\ST@pageleft -\ST@toadd
    \global\ST@toadd=\z@}}
\else
%    \end{macrocode}
% 以下是较新版本。但在 \texttt{supertabular} 4.2a 及更高版本中，\cs{ST@lineht} 更名为 \cs{ST@rowht}。
%
%    \begin{macrocode}
\def\ST@tabularcr{%
  {\ifnum0=`}\fi
  \ST@save@lineno
  \@ifstar{\MRST@xtabularcr}{\ST@xtabularcr}}
\ifx\ST@rowht\undefined
\def\MRST@cr{%
  \noalign{%
    \ifnum\ST@pboxht<\ST@lineht
      \global\advance\ST@pageleft -\ST@lineht
      \global\ST@prevht\ST@lineht
    \else
     \ST@trace@cr\thr@@{Added par box with height \the\ST@pboxht}%
      \global\advance\ST@pageleft -\ST@pboxht
      \global\advance\ST@pageleft -0.1\ST@pboxht
      \global\ST@prevht\ST@pboxht
      \global\ST@pboxht\z@
    \fi
    \global\advance\ST@pageleft -\ST@toadd
    \global\ST@toadd=\z@
    \ST@trace@cr\thr@@{Space left for tabular: \the\ST@pageleft}}}
\else
%    \end{macrocode}
% 下面是 4.2a 之后的版本。
%    \begin{macrocode}
\def\MRST@cr{%
  \noalign{%
    \ifnum\ST@pboxht<\ST@rowht
      \global\advance\ST@pageleft -\ST@rowht
      \global\ST@prevht\ST@rowht
    \else
     \ST@trace@cr\thr@@{Added par box with height \the\ST@pboxht}%
      \global\advance\ST@pageleft -\ST@pboxht
      \global\advance\ST@pageleft -0.1\ST@pboxht
      \global\ST@prevht\ST@pboxht
      \global\ST@pboxht\z@
    \fi
    \global\advance\ST@pageleft -\ST@toadd
    \global\ST@toadd=\z@
    \ST@trace@cr\thr@@{Space left for tabular: \the\ST@pageleft}}}
\fi
\fi
%    \end{macrocode}
% \end{macro}
% \begin{macro}{\MRST@xtabularc}
% \begin{macro}{\MRST@argtabularc}
% \begin{macro}{\MRST@xargtabularc}
% \begin{macro}{\MRST@yargtabularc}
% 以下是 \Package{supertabular} 中相应宏的副本，但它们调用 \cs{MRST@cr} 而非 \cs{ST@cr}。
%
%    \begin{macrocode}
\def\MRST@xtabularcr{%
  \@ifnextchar[%]
    {\MRST@argtabularcr}%
    {\ifnum0=`{\fi}\cr\MRST@cr}}
\def\MRST@argtabularcr[#1]{%
  \ifnum0=`{\fi}%
  \ifdim #1>\z@
    \unskip\MRST@xargarraycr{#1}
  \else
    \MRST@yargarraycr{#1}%
  \fi}
\def\MRST@xargarraycr#1{%
  \@tempdima #1\advance\@tempdima \dp \@arstrutbox
  \vrule \@height\z@ \@depth\@tempdima \@width\z@ \cr
  \noalign{\global\ST@toadd=#1}\MRST@cr}
\def\MRST@yargarraycr#1{%
  \cr\noalign{\vskip #1\global\MRST@toadd=#1}\MRST@cr}
}
%    \end{macrocode}
% \end{macro}
% \end{macro}
% \end{macro}
% \end{macro}
% \begin{macro}{\STneed}
% 该宏可在 \env{supertabular} 中使用，用以说明一个跨行条目需要多少空间。参见第~\ref{sec:use-with-supertabular} 节。
%
%    \begin{macrocode}
\def\STneed#1{\ifdim\ST@pageleft<#1\ST@newpage\ST@next\fi}
}
%    \end{macrocode}
% \end{macro}
%    \begin{macrocode}
\ProcessOptions
%    \end{macrocode}
%
% \cs{multirow@colwidth} 是一个长度寄存器，用于实现 \meta{width} 的“\texttt{=}”形式。
%
% \changes{v1.9}{2016/09/15}{为 \cs{multirow} 的 \meta{width} 参数实现等号形式的选项。}
% \begin{macro}{\multirow@colwidth}
% \cs{multirow@colwidth} 是一个长度寄存器，用于实现 \meta{width} 的“\texttt{=}”形式。
%    \begin{macrocode}
\newlength{\multirow@colwidth}
%    \end{macrocode}
% \end{macro}
%
% \changes{v1.9}{2016/09/15}{为 \Package{multirow} 提供独立的临时
% 寄存器，以便把盒子高度安全地传递给 \Package{bigdelim}。}
% \changes{v2.2}{2016/11/25}{移除 \cs{multirow@cnta}。}
% \begin{macro}{\multirow@cntb}
% \begin{macro}{\multirow@dima}
%   定义两个计数器和一个长度，供 \cs{multirow} 内部使用。
%    \begin{macrocode}
\newcount\multirow@cntb
\newlength\multirow@dima
%    \end{macrocode}
% \end{macro}
% \end{macro}
% \begin{macro}{\multirow@setcolwidth}
% 该宏为宽度取“\texttt{=}”的条目计算 \cs{multirow@colwidth}。我们通过检查 \cs{TY@final} 是否定义来判断是否处于 \env{tabulary} 环境中；若否，则 \cs{multirow@colwidth} = \cs{hsize}。\env{tabulary} 会进行两遍排版：第一遍把 \cs{multirow@colwidth} 设为文本在 LR 模式下的尺寸（换行符替换为空格），以使 \env{tabulary} 给出足够空间；第二遍（以 \cs{TY@box} = \cs{TY@box@v} 为标志）则采用 \env{tabulary} 在 \cs{hsize} 中给出的值。该算法并不完美，但在多数情况下足够好用。
%
%
%
%
%
%
%
%
%
% \changes{v2.7}{2021/01/29}{将 \cs{multirow@setcolwidth} 改为 \cs{long} 定义，以支持多段文本}
%
%    \begin{macrocode}
\long\def\multirow@setcolwidth#1{%
    \ifx\TY@final\multirow@undefined \multirow@colwidth=\hsize
    \else
      \ifx\TY@box\TY@box@v\multirow@colwidth=\hsize
      \else \setbox0\hbox
        {\let\\\space\let\newline\space #1}\multirow@colwidth=\wd0
      \fi
    \fi}
%    \end{macrocode}
% \end{macro}
%
% \begin{macro}{\multirowsetup}
% \cs{multirowsetup} 在每个 \cs{multirow} 开始处执行。
%    \begin{macrocode}
\newcommand\multirowsetup{\raggedright}
%    \end{macrocode}
% \end{macro}
%
% \begin{macro}{\multirow@vbox}
% 生成 \cs{vbox}。参数：\\
% \texttt{\#1} = \meta{vpos}，\texttt{\#2} = 初始化代码（例如设置 \cs{parbox} 的宽度），\texttt{\#3} = 盒子内容。\\
%
% 根据 \meta{vpos} 参数，盒子将顶部对齐、垂直居中或底部对齐；这通过在相应位置插入 \cs{vfill} 实现。\\
%
%
% 注意：\cs{relax} 用于防止 \meta{vpos} 参数为空。
%
%    \begin{macrocode}
\long\def\multirow@vbox#1#2#3{\setbox0\vtop to \multirow@dima{#2%
    \if #1t\relax\else\vfill\fi
    \multirowsetup #3\if #1b\relax\else\vfill\fi}}
%    \end{macrocode}
% \end{macro}
%
% \begin{macro}{\multirow}
% 生成一个跨越表格多行的条目。\\
% 首先收集所有参数，并用默认值补齐缺失的可选参数。
% \changes{v1.8}{2016/09/14}{增加第一个可选参数 \meta{vpos}。}
%    \begin{macrocode}
%% \multirow [vpos] {nrows} [bigstruts] {width} [vmove] {text}
\newcommand\multirow[2][c]{\@multirow[#1]{#2}}
\def\@multirow[#1]#2{\@ifnextchar[{\@@multirow[#1]#2}{\@@multirow[#1]#2[0]}}
\def\@@multirow[#1]#2[#3]#4{\@ifnextchar[{\@xmultirow[#1]{#2}[#3]{#4}}%
                                {\@xmultirow[#1]{#2}[#3]{#4}[0pt]}}
%    \end{macrocode}
% \changes{v1.9a}{2016/09/23}{为 \meta{bigstruts} 参数增加可选前缀。}
% \begin{macro}{\multirow@piii}
% \begin{macro}{\ifmultirow@prefixt}
% \begin{macro}{\multirow@prefixttrue}
% \begin{macro}{\multirow@prefixtfalse}
% \begin{macro}{\ifmultirow@prefixb}
% \begin{macro}{\multirow@prefixbtrue}
% \begin{macro}{\multirow@prefixbfalse}
% 该宏从 \meta{bigstruts} 参数中分离出 \texttt{t}、\texttt{b} 或 \texttt{tb} 前缀，并把数值赋给 \cs{multirow@cntb}。前缀记录在两个布尔开关 \cs{ifmultirow@prefixt} 与 \cs{ifmultirow@prefixb} 中。
%
%
%
%    \begin{macrocode}
\newif\ifmultirow@prefixt
\newif\ifmultirow@prefixb
\def\multirow@piii#1#2#3\end{\multirow@prefixtfalse\multirow@prefixbfalse
  \if t#1\multirow@prefixttrue
    \if b#2\multirow@prefixbtrue \multirow@cntb=#3%
    \else \multirow@cntb=#2#3%
    \fi
  \else
    \if b#1\multirow@prefixbtrue \multirow@cntb=#2#3%
    \else \multirow@cntb=#1#2#3%
    \fi
  \fi}
%    \end{macrocode}
% \end{macro}
% \end{macro}
% \end{macro}
% \end{macro}
% \end{macro}
% \end{macro}
% \end{macro}
%
% 这是真正的核心宏。先拆分 \meta{bigstruts} 参数，然后计算跨行盒子的高度。
%
% \changes{v2.2}{2016/11/25}{支持 \meta{nrows} 取小数值。}
% 由于 \meta{nrows}（\verb+#2+）可能为小数，不能用 \cs{ifnum} 判断正负，因此在数字后加上单位（\texttt{pt}）改用 \cs{ifdim}。
%
%
% \changes{v2.4}{2019/01/01}{支持 \meta{width} 与 \meta{vmove} 使用与 \Package{calc} 兼容的表达式。}
%
% \changes{v2.7}{2021/01/29}{将 \cs{@xmultirow} 改为 \cs{long} 定义，以支持多段文本}
%
%    \begin{macrocode}
\long\def\@xmultirow[#1]#2[#3]#4[#5]#6{%
  \expandafter\multirow@piii#3\relax\end%
  \setlength\multirow@dima{#2\ht\@arstrutbox}%
  \addtolength\multirow@dima{#2\dp\@arstrutbox}%
  \ifdim#2pt<\z@\setlength\multirow@dima{-\multirow@dima}\fi
  \addtolength\multirow@dima{\multirow@cntb\bigstrutjot}%
%    \end{macrocode}
% \changes{v1.9a}{2016/09/23}{重做 \cs{vbox} 的计算与定位。}
% 文本通过调用 \cs{multirow@vbox} 排入 \cs{vbox}。\\
% 若 \meta{width} 参数为 \texttt{*}，则只把文本放入 \cs{vbox}。
%    \begin{macrocode}
  \if*#4\multirow@vbox{#1}{}{\hbox{\strut#6\strut}}%
%    \end{macrocode}
% 否则把文本放入 \cs{vbox} 内的 \cs{parbox} 中。\\
% 若 \meta{width} 参数取“\texttt{=}”，则先计算 \cs{multirow@colwidth}，并将其作为 \cs{parbox} 的宽度。
%
%    \begin{macrocode}
  \else \if=#4\multirow@setcolwidth{#6}%
    \multirow@vbox{#1}{\setlength\hsize{\multirow@colwidth}\@parboxrestore}{\strut#6\strut\par}%
%    \end{macrocode}
% 否则把给定参数用作 \cs{parbox} 的宽度。
%    \begin{macrocode}
  \else \multirow@vbox{#1}{\setlength\hsize{#4}\@parboxrestore}{\strut#6\strut\par}%
  \fi \fi
%    \end{macrocode}
% 接下来正确定位 \cs{vbox}。细节见附录；平移量计算的概览见第~\ref{sec:overview} 节。
%
%
%
% 当 \meta{nrows}${}> 0$ 时：\\
% 若 \meta{vpos} = \texttt{[t]}，盒子已定位正确（基线位于该行基线上）。但稍后会以盒子顶部（而非基线）作为参考点，因此取盒子高度 h 作为平移量。见图~\ref{fig:top-t}。\\
%
%
%
%
% 若 \meta{vpos} = \texttt{[c]}，上移 h1（见图~\ref{fig:top-c}），其中 h1 = \cs{ht}\cs{@arstrutbox} +（\cs{bigstrutjot} \cs{ifmultirow@prefixt}）。\\
%
% 若 \meta{vpos} = \texttt{[b]}，上移 h1 + h2（见图~\ref{fig:top-b}），其中 h2 = \cs{dp}\cs{@arstrutbox} +（\cs{bigstrutjot} \cs{ifmultirow@prefixb}）。\\
%
% 所需平移量在 \cs{multirow@dima} 中计算。
%    \begin{macrocode}
  \ifdim#2pt>\z@
    \if#1t\relax\setlength\multirow@dima{\ht0}\else
      \setlength\multirow@dima{\ht\@arstrutbox}%
      \ifmultirow@prefixt \addtolength\multirow@dima{\bigstrutjot}\fi
      \if#1b\relax \addtolength\multirow@dima{\dp\@arstrutbox}%
        \ifmultirow@prefixb \addtolength\multirow@dima{\bigstrutjot}\fi
      \fi
    \fi
%    \end{macrocode}
% 当 \meta{nrows}${}< 0$ 时：\\
% 若 \meta{vpos} = \texttt{[t]}，把盒子上移 $\textrm{H} - \textrm{h1} -\textrm{h2} + \textrm{h}$。见图~\ref{fig:bot-t}。\\
%
% 若 \meta{vpos} = \texttt{[c]}，把盒子上移 $\textrm{H} - \textrm{h2}$。见图~\ref{fig:bot-c}。\\
% 若 \meta{vpos} = \texttt{[b]}，把盒子上移 H。见图~\ref{fig:bot-b}。\\
% H 为 \cs{multirow@dima} 的当前值。
%    \begin{macrocode}
  \else
    \if#1b\relax\else
      \addtolength\multirow@dima{-\dp\@arstrutbox}%
      \ifmultirow@prefixb \addtolength\multirow@dima{-\bigstrutjot}\fi
      \if#1t\relax\addtolength\multirow@dima{-\ht\@arstrutbox}%
        \ifmultirow@prefixt \addtolength\multirow@dima{-\bigstrutjot}\fi
        \addtolength\multirow@dima{\ht0}%
      \fi
    \fi
  \fi
%    \end{macrocode}
% \changes{v2.1}{2016/10/11}{把最终 \cs{vbox} 的深度设为 0，以防止过高的跨行行把后续行向下推。}
%
% 最后加上 \meta{vmove} 参数（\verb+#5+），进入水平模式。然后在盒子上方放置 \cs{vskip} 使其上移，再加入输出。由于该 \cs{vskip}，所得盒子高度为 0；我们把 \cs{vbox} 的深度也设为 0，以免影响当前行的深度。\\
%
%
%
%
% 若 \cs{multirowdebug} 为真，则显示该盒子。
%    \begin{macrocode}
  \addtolength\multirow@dima{#5}%
  \leavevmode
  \setbox0\vtop{\vskip-\multirow@dima\box0\vss}\dp0=\z@
  \ifmultirowdebug{\showboxdepth=5 \showboxbreadth=10 \showbox0}\fi
  \box0
}
%    \end{macrocode}
% \end{macro}
% \begin{macro}{\bigstrutjot}
% 若尚未定义，则定义 \cs{bigstrutjot}。
%    \begin{macrocode}
\@ifundefined{bigstrutjot}{\newdimen\bigstrutjot \bigstrutjot=\jot}{}
%    \end{macrocode}
% \end{macro}
% \iffalse
%</multirow>
% \fi
%
% \subsection{\Package{bigstrut} 宏包}
% \label{sec:pack-bigstrut}
%
% \iffalse
%<*bigstrut>
% \fi
% \changes{bigstrut v1.0}{1994/05/31}{初始版本}
% \begin{macro}{\bigstrutjot}
% 这是一个长度，默认值为 2pt。可用 \cs{setlength} 命令修改。
%
%    \begin{macrocode}
\@ifundefined{bigstrutjot}{\newdimen\bigstrutjot}{}\bigstrutjot=2pt
%    \end{macrocode}
% \end{macro}
%
% \changes{bigstrut v2.4}{2018/12/30}{在开头增加 \cs{leavevmode}
% 以强制进入水平模式}
% \begin{macro}{\bigstrut}
% 该宏插入一个支柱。根据可选参数的不同，它会在标准 \env{array}/\env{tabular} 支柱的上方和/或下方延伸。
%
%    \begin{macrocode}
\newcommand\bigstrut[1][x]{%
  \leavevmode\unskip\@tempdima=\ht\@arstrutbox \@tempdimb=\dp\@arstrutbox
  \ifx #1b\relax \else \advance\@tempdima by \bigstrutjot\fi
  \ifx #1t\relax \else \advance\@tempdimb by \bigstrutjot\fi
  \hbox{\vrule \@height\@tempdima \@depth\@tempdimb \@width\z@}\ignorespaces}
%    \end{macrocode}
% \end{macro}
%
% \iffalse
%</bigstrut>
% \fi
%
% \subsection{\Package{bigdelim} 宏包}
% \label{sec:pack-bigdelim}
% \iffalse
%<*bigdelim>
% \fi
% \changes{bigdelim v0.0}{1994/10/02}{\texttt{bigbrace.sty}，作者 \O ystein Bache}
% \changes{bigdelim v1.0}{1999/11/05}{初始版本 \texttt{bigdelim.sty}}
%    \begin{macrocode}
\RequirePackage{multirow}
%    \end{macrocode}
% \begin{macro}{\ldelim}
% 该宏排版一个左侧定界符。它以合适的参数调用 \cs{multirow}。
% 定界符的尺寸通过在其旁放置一个高度合适、宽度为零的 \cs{vbox} 来确定；该高度即 \cs{multirow} 已在 \cs{multirow@dima} 中算出的值。
%
% 该计算使用 \cs{@arstrutbox} 的尺寸，而后者由 \texttt{tabular} 或 \texttt{array} 环境设置；若尚未设置，则将其初始化为默认值。
%
%
%
% \changes{bigdelim v2.3}{2018/08/03}{将 \cs{textrm} 替换为 \cs{textnormal}}
% \changes{bigdelim v2.6}{2020/12/26}{若 \cs{@arstrutbox} 未定义则初始化它}
% \changes{bigdelim v2.8}{2021/03/15}{增加可选参数 \meta{vmove}}
%    \begin{macrocode}
\newcommand\ldelim[2]{\@ifnextchar[{\@ldelim{#1}{#2}}{\@ldelim{#1}{#2}[0pt]}}
\def\@ldelim#1#2[#3]#4{\@ifnextchar[{\@@ldelim{#1}{#2}{#3}{#4}}{\@@ldelim{#1}{#2}{#3}{#4}[\null]}}
\def\@@ldelim#1#2#3#4[#5]%
 {\ifvoid\@arstrutbox\setbox\@arstrutbox\hbox{\strut}\fi
  \multirow{#2}{#4}[#3]{%
   \ensuremath
     {\left.\vcenter{\hsize=0pt\vrule height \multirow@dima width 0pt}%
     \textnormal{#5}\right#1}}}
%    \end{macrocode}
% \end{macro}
%
% \begin{macro}{\rdelim}
% 该宏排版一个右侧定界符。它以合适的参数调用 \cs{multirow}，做法与 \cs{ldelim} 类似。
%
% \changes{bigdelim v2.3}{2018/08/03}{将 \cs{textrm} 替换为 \cs{textnormal}}
% \changes{bigdelim v2.6}{2020/12/26}{若 \cs{@arstrutbox} 未定义则初始化它}
% \changes{bigdelim v2.8}{2021/03/15}{增加可选参数 \meta{vmove}}
%    \begin{macrocode}
\newcommand\rdelim[2]{\@ifnextchar[{\@rdelim{#1}{#2}}{\@rdelim{#1}{#2}[0pt]}}
\def\@rdelim#1#2[#3]#4{\@ifnextchar[{\@@rdelim{#1}{#2}{#3}{#4}}{\@@rdelim{#1}{#2}{#3}{#4}[\null]}}
\def\@@rdelim#1#2#3#4[#5]%
 {\ifvoid\@arstrutbox\setbox\@arstrutbox\hbox{\strut}\fi
  \multirow{#2}{#4}[#3]{%
   \ensuremath
     {\left#1\vcenter{\hsize=0pt\vrule height \multirow@dima width 0pt}%
       \textnormal{#5}\right.}}}
%    \end{macrocode}
% \end{macro}
% \iffalse
%</bigdelim>
% \fi
%
% \appendix
%
% \section{附录}
% \label{sec:appendix}
%
% \index{vbox=\verb!*+\vbox+|usage}
%
% 本节说明 \cs{multirow} 中 \cs{vbox} 的定位方式。
% 定位方式取决于 \meta{nrows}、\meta{vpos}、
% \meta{bigstruts} 与 \meta{vmove} 参数。盒子用 \cs{vtop} 构造。
% \cs{vtop} 的算法见 \emph{The \TeX{}book} 第 81 页。
%
% 每种情形都用一幅图说明。图中左列
% 表示跨行条目所在表格的上下文，即 \meta{nrows} 行；右列是待插入的跨行盒子。基线是内容最初所处的自然位置，之后再上移到目标位置。
%
%
%
%
%
% H 为算出的盒子高度：$\meta{nrows} \times
% \textrm{一行的自然高度}  + \meta{bigstruts} \times
% \mathrm{\cs{bigstrutjot}} $。
%
% {顶支柱}：若第一行顶部有 \cs{bigstrut}（由 \meta{bigstruts} 参数中的 \texttt{t} 前缀表示），则等于 \cs{bigstrutjot}，否则为 0。
%
%
%
% {底支柱}：若最后一行底部有 \cs{bigstrut}（由 \meta{bigstruts} 参数中的 \texttt{b} 前缀表示），则等于 \cs{bigstrutjot}，否则为 0。
%
%
%
% h1 = $\textrm{表格行的高度} + \textrm{顶支柱}$
%
% h2 = $\textrm{表格行的深度} + \textrm{底支柱}$
%
% 注意：以下描述给出的是不计 \meta{vmove} 时盒子的垂直平移量；无论哪种情形，若给出了 \meta{vmove}，都需再加上它。
%
%
%
% \subsection{情形 \meta{nrows}${}> 0$}
%
% \subsubsection*{\meta{vpos}  = [t]}
%
% 此时 \cs{vbox} 中先是文本，后接一个 \cs{vfill}。
% 这样的 \cs{vbox} 高度等于文本首行的高度（h）。H = 盒子高度 + 盒子深度，这意味着盒子已定位正确。但稍后我们会把盒子放入另一个 \cs{vbox}，其顶部有一个 \cs{vskip}，这会使盒子顶部成为参考点。因此还需再上移距离 h（它很可能与表格行的高度不同），总平移量即为 h。见图~\ref{fig:top-t}。
%
%
%
%
%
%
%
%
% 另一种做法是在此情形省略 \cs{vskip}，从而保持基线不动；但这会使代码不对称，而且在存在非零 \meta{vmove} 时无法工作。因此我们选择在此把平移量设为 h。
%
%
%
%
% \begin{figure}[htpb]
%   \centering
% \begin{tikzpicture}
%   \draw (2, -0.5) rectangle (5,4.5);
%   \foreach \y in {0,...,4} \draw (2, \y) -- (5, \y);
%   \draw (0, 4.2) node (topstrut) {顶支柱};
%   \draw[->] (topstrut) -- (2, 4.2);
%   \draw (0, -0.25) node (botstrut) {底支柱};
%   \draw[->] (botstrut) -- (2, -0.25);
%   \draw[<->] (5.3, -0.5) -- (5.3, 4.5);
%   \draw (5.5, 2) node {H};
%   \draw (0, 3.3) node (baseline) {基线};
%   \draw[dashed] (1.7, 3.3) -- (5.9, 3.3);
%   \draw[<->] (5.5, 3.3) -- (5.5, 4.5);
%   \draw (5.9, 3.9) node {h1};
%
%   \draw (6.5, 4) rectangle (9.5, -1);
%   \draw[<->] (9.7, -1) -- (9.7, 4);
%   \draw (10, 1.5) node {H};
%   \draw[<->] (10, 3.3) -- (10, 4);
%   \draw (10.3, 3.7) node {h};
%   \draw[fill=lightgray] (6.6, 1) rectangle (9.4, 3.9);
%   \draw (8, 0)  node[align=center] {\cs{vfill}};
%   \draw (8, 2) node[align=center] {文本};
%   \draw[draw=gray] (6.6, 3) -- (9.4, 3);
%   \draw[dashed] (6.2, 3.3) -- (9.8, 3.3);
%   \draw[->, line width=4pt, draw=gray] (8, 3.3) -- (8, 4);
% \end{tikzpicture}
% \caption{情形 \meta{nrows}${}> 0$，\meta{vpos} = [t]}
% \label{fig:top-t}
% \end{figure}
%
% \subsubsection*{\meta{vpos}  = [c]}
%
% 此时 \cs{vbox} 中依次为一个 \cs{vfill}、文本和另一个 \cs{vfill}。
% 这样的 \cs{vbox} 高度为 0，即盒子顶部位于基线上。由于两个盒子尺寸相同（H），把 \cs{vbox} 上移 h1 即可对齐。见图~\ref{fig:top-c}。
%
%
%
% \begin{figure}[htpb]
%   \centering
%\begin{tikzpicture}
%   \draw (2, -0.5) rectangle (5,4.5);
%   \foreach \y in {0,...,4} \draw (2, \y) -- (5, \y);
%   \draw (0, 4.2) node (topstrut) {顶支柱};
%   \draw[->] (topstrut) -- (2, 4.2);
%   \draw (0, -0.25) node (botstrut) {底支柱};
%   \draw[->] (botstrut) -- (2, -0.25);
%   \draw[<->] (5.3, -0.5) -- (5.3, 4.5);
%   \draw (5.5, 2) node {H};
%   \draw (0, 3.3) node (baseline) {基线};
%   \draw[dashed] (1.7, 3.3) -- (5.9, 3.3);
%   \draw[<->] (5.5, 3.3) -- (5.5, 4.5);
%   \draw (5.9, 3.9) node {h1};
%
%   \draw (6.5, 3.3) rectangle (9.5, -1.7);
%   \draw[<->] (9.7, -1.7) -- (9.7, 3.3);
%   \draw (10, 0.8) node {H};
%   \draw[fill=lightgray] (6.6, -0.5) rectangle (9.4, 2.1);
%   \draw (8, 2.7)  node[align=center] {\cs{vfill}};
%   \draw (8, 0.8) node[align=center] {文本};
%   \draw (8, -1) node[align=center] {\cs{vfill}};
%   \draw[->, line width=4pt, draw=gray] (8, 3.3) -- (8, 4.5);
%   \draw[draw=gray] (6.5, 4.5) -- (9.5, 4.5);
% \end{tikzpicture}
% \caption{情形 \meta{nrows}${}> 0$，\meta{vpos} = [c]}
% \label{fig:top-c}
% \end{figure}
%
% \subsubsection*{\meta{vpos}  = [b]}
%
% 此时 \cs{vbox} 中先是一个 \cs{vfill}，后接文本。
% 由于以文本结尾，盒子还会多出一段等于文本末行深度的深度。
%
% 这样的 \cs{vbox} 高度为 0，即盒子顶部位于基线上，但其深度为 H 加上上述深度；换言之，末行文本的基线位于顶部以下 H 处。
%
%
%
% 由于 \meta{vpos} = \texttt{[b]}，我们希望末行文本的基线平移到表格最后一行的基线，平移量为 h1 + h2。见图~\ref{fig:top-b}。
%
%
%
% \begin{figure}[htpb]
%   \centering
%\begin{tikzpicture}
%   \draw (2, -0.5) rectangle (5,4.5);
%   \foreach \y in {0,...,4} \draw (2, \y) -- (5, \y);
%   \draw (0, 4.2) node (topstrut) {顶支柱};
%   \draw[->] (topstrut) -- (2, 4.2);
%   \draw (0, -0.25) node (botstrut) {底支柱};
%   \draw[->] (botstrut) -- (2, -0.25);
%   \draw[<->] (5.3, -0.5) -- (5.3, 4.5);
%   \draw (5.5, 2) node {H};
%   \draw (0, 3.3) node (baseline) {基线};
%   \draw[dashed] (1.7, 3.3) -- (5.9, 3.3);
%   \draw[<->] (5.5, 3.3) -- (5.5, 4.5);
%   \draw (5.9, 3.9) node {h1};
%   \draw[<->] (5.5, -0.5) -- (5.5, 0.3);
%   \draw (5.9, -0.1) node {h2};
%
%   \draw (6.5, 3.3) rectangle (9.5, -2.1);
%   \draw[<->] (9.7, -1.7) -- (9.7, 3.3);
%   \draw (10, 0.8) node {H};
%   \draw[fill=lightgray] (6.6, -2) rectangle (9.4, 1);
%   \draw (8, 2)  node[align=center] {{\cs{vfill}}};
%   \draw (8, -0.5) node[align=center] {文本};
%   \draw (0, 0.3) node (baseline) {基线};
%   \draw[dashed] (6.3, -1.7) -- (10, -1.7);
%   \draw (3.5, -1.7) node (baseline) {末行文本的基线};
%   \draw[dashed] (1.7, 0.3) -- (10, 0.3);
%   \draw[->, line width=4pt, draw=gray] (8.5, -1.7) -- (8.5, 0.3);
% \end{tikzpicture}
% \caption{情形 \meta{nrows}${}> 0$，\meta{vpos} = [b]}
% \label{fig:top-b}
% \end{figure}
%
% \subsection{情形 \meta{nrows}${}<0$}
%
% 当跨行条目位于跨行块的最后一行时，\meta{nrows}${}<0$。
%
%
% \subsubsection*{\meta{vpos}  = [t]}
%
% 此时 \cs{vbox} 中先是文本，后接一个 \cs{vfill}。
% 这样的 \cs{vbox} 高度等于文本首行的高度。基线与最后一行的基线对齐。
%
% 由于 \meta{vpos} = [t]，我们希望它与第一行的基线对齐，因此需上移 $\textrm{H} - \textrm{h1} - \textrm{h2}$。但稍后盒子高度会被置为 0，所以还必须加上当前高度 h，总平移量为 $\textrm{H} - \textrm{h1} - \textrm{h2} + \textrm{h}$。见图~\ref{fig:bot-t}。
%
%
%
%
%
%
% \begin{figure}[htpb]
%   \centering
% \begin{tikzpicture}
%   \draw (2, -0.5) rectangle (5,4.5);
%   \foreach \y in {0,...,4} \draw (2, \y) -- (5, \y);
%   \draw (0, 4.2) node (topstrut) {顶支柱};
%   \draw[->] (topstrut) -- (2, 4.2);
%   \draw (0, -0.25) node (botstrut) {底支柱};
%   \draw[->] (botstrut) -- (2, -0.25);
%   \draw[<->] (5.3, -0.5) -- (5.3, 4.5);
%   \draw (5.5, 2) node {H};
%   \draw (0, 3.3) node (baseline) {基线};
%   \draw[dashed] (1.7, 3.3) -- (9, 3.3);
%   \draw[<->] (5.5, 3.3) -- (5.5, 4.5);
%   \draw (5.9, 3.9) node {h1};
%   \draw[<->] (5.5, -0.5) -- (5.5, 0.3);
%   \draw (5.9, -0.1) node {h2};
%   \draw (0, 0.3) node (baseline) {基线};
%   \draw[dashed] (1.7, 0.3) -- (5.9, 0.3);
%
%   \draw (6.5, 1) rectangle (9.5, -4);
%   \draw[<->] (9.7, -4) -- (9.7, 1);
%   \draw (10, -1.5) node {H};
%   \draw[<->] (10, 0.3) -- (10, 1);
%   \draw (10.3, 0.7) node {h};
%   \draw[fill=lightgray] (6.6, -2) rectangle (9.4, 0.9);
%   \draw (8, -3)  node[align=center] {\cs{vfill}};
%   \draw (8, -1) node[align=center] {文本};
%   \draw[draw=gray] (6.6, 0) -- (9.4, 0);
%   \draw[dashed] (6.2, 0.3) -- (10.3, 0.3);
%   \draw[->, line width=4pt, draw=gray] (8, 0.3) -- (8, 3.3);
% \end{tikzpicture}
% \caption{情形 \meta{nrows}${} < 0$，\meta{vpos} = [t]}
% \label{fig:bot-t}
% \end{figure}
%
% \subsubsection*{\meta{vpos}  = [c]}
%
% 此时 \cs{vbox} 中依次为一个 \cs{vfill}、文本和另一个 \cs{vfill}。
% 这样的 \cs{vbox} 高度为 0，即盒子顶部位于基线上。由于两个盒子尺寸相同（H），把 \cs{vbox} 上移 $\textrm{H} - \textrm{h2}$ 即可对齐。见图~\ref{fig:bot-c}。
%
%
%
%
% \begin{figure}[htpb]
% \begin{tikzpicture}
%   \draw (2, -0.5) rectangle (5,4.5);
%   \foreach \y in {0,...,4} \draw (2, \y) -- (5, \y);
%   \draw (0, 4.2) node (topstrut) {顶支柱};
%   \draw[->] (topstrut) -- (2, 4.2);
%   \draw (0, -0.25) node (botstrut) {底支柱};
%   \draw[->] (botstrut) -- (2, -0.25);
%   \draw[<->] (5.3, -0.5) -- (5.3, 4.5);
%   \draw (5.5, 2) node {H};
%   \draw (0, 3.3) node {基线};
%   \draw[dashed] (1.7, 3.3) -- (5.8, 3.3);
%   \draw[<->] (5.5, 3.3) -- (5.5, 4.5);
%   \draw (5.9, 3.9) node {h1};
%   \draw[<->] (5.5, -0.5) -- (5.5, 0.3);
%   \draw (5.9, -0.1) node {h2};
%   \draw (0, 0.3) node {基线};
%   \draw[dashed] (1.7, 0.3) -- (5.9, 0.3);
%
%   \draw (6.5, 0.3) rectangle (9.5, -4.7);
%   \draw[<->] (9.7, -4.7) -- (9.7, 0.3);
%   \draw (10, -3.2) node {H};
%   \draw[fill=lightgray] (6.6, -3.5) rectangle (9.4, -0.9);
%   \draw (8, -0.3)  node[align=center] {\cs{vfill}};
%   \draw (8, -2.2) node[align=center] {文本};
%   \draw (8, -4) node[align=center] {\cs{vfill}};
%   \draw[->, line width=4pt, draw=gray] (8, 0.3) -- (8, 4.5);
%   \draw[draw=gray] (6.5, 4.5) -- (9.5, 4.5);
% \end{tikzpicture}
% \caption{情形 \meta{nrows}${} < 0$，\meta{vpos} = [c]}
% \label{fig:bot-c}
% \end{figure}
%
% \subsubsection*{\meta{vpos}  = [b]}
%
% \cs{vbox} 中先是一个 \cs{vfill}，后接文本。
% 由于以文本结尾，盒子还会多出一段等于文本末行深度的深度。
%
% 这样的 \cs{vbox} 高度为 0，即盒子顶部位于基线上，但其深度为 H 加上上述深度；换言之，末行文本的基线位于顶部以下 H 处。
%
%
%
% 由于 \meta{vpos} = \texttt{[b]}，我们希望末行文本的基线平移到表格最后一行的基线，平移量为 H。见图~\ref{fig:bot-b}。
%
%
%
% \begin{figure}[htpb]
% \begin{tikzpicture}
%   \draw (2, -0.5) rectangle (5,4.5);
%   \foreach \y in {0,...,4} \draw (2, \y) -- (5, \y);
%   \draw (0, 4.2) node (topstrut) {顶支柱};
%   \draw[->] (topstrut) -- (2, 4.2);
%   \draw (0, -0.25) node (botstrut) {底支柱};
%   \draw[->] (botstrut) -- (2, -0.25);
%   \draw[<->] (5.3, -0.5) -- (5.3, 4.5);
%   \draw (5.5, 2) node {H};
%   \draw (0, 3.3) node {基线};
%   \draw[dashed] (1.7, 3.3) -- (5.8, 3.3);
%   \draw[<->] (5.5, 3.3) -- (5.5, 4.5);
%   \draw (5.9, 3.9) node {h1};
%   \draw[<->] (5.5, -0.5) -- (5.5, 0.3);
%   \draw (5.9, -0.1) node {h2};
%   \draw (0, 0.3) node {基线};
%   \draw[dashed] (1.7, 0.3) -- (5.9, 0.3);
%
%   \draw (6.5, 0.3) rectangle (9.5, -5.1);
%   \draw[<->] (9.7, -4.7) -- (9.7, 0.3);
%   \draw (10, -3.2) node {H};
%   \draw[fill=lightgray] (6.6, -5) rectangle (9.4, -2);
%   \draw (8, -1)  node[align=center] {{\cs{vfill}}};
%   \draw (8, -3.5) node[align=center] {文本};
%   \draw (0, 0.3) node (baseline) {基线};
%   \draw[dashed] (1.7, 0.3) -- (10, 0.3);
%   \draw[dashed] (6.3, -4.7) -- (10, -4.7);
%   \draw (3.5, -4.7) node {末行文本的基线};
%   \draw[->, line width=4pt, draw=gray] (8.5, -4.7) -- (8.5, 0.3);
% \end{tikzpicture}
% \caption{情形 \meta{nrows}${} < 0$，\meta{vpos} = [b]}
% \label{fig:bot-b}
% \end{figure}
%
% \subsection{概览}
% \label{sec:overview}
%
% \begin{tabular}{| l | l | l |}
%     \hline
%   \meta{vpos} & \meta{nrows} > 0  & \meta{nrows} < 0 \\
%   \hline
%   \texttt{[t]}  & h            &  $\textrm{H} - \textrm{h1} -\textrm{h2} + \textrm{h}$  \\
%   \texttt{[c]}  & h1          &  $\textrm{H} - \textrm{h2}$ \\
%   \texttt{[b]}  & h1 + h2 &   H \\
%   \hline
%                    &  x          &   $\textrm{H} - \textrm{h1} -\textrm{h2} + x$ \\
% \hline
%   \end{tabular}
%
% \Finale
\endinput
